% Agent contribution: /root/archive_questions.
% AQ identifiers denote findings of this run, not claims of priority.
\section{Archive questions: exact statements and mathematical scope}

\subsection{AQ-FINITE: finite Frobenius recognition of abelian Artin--Tate motives}
\paragraph{Statement.}
Adopt the Dirichlet GRH theorem supplied in \texttt{ALPOGE\_Main.lean},
\texttt{LFunction\_eq\_zero\_re}. Let $N\geq1$, put
\[
 Q=\operatorname{lcm}(N,4),\qquad B_Q=(\varphi(Q)\log Q)^2,
\]
and let $M,M'$ be finite direct sums of Tate twists of abelian Artin
motives over $\mathbb Q$, whose Artin representations factor through
$G_N=(\mathbb Z/N\mathbb Z)^\times$. Work initially with a coefficient
number field $E\subset\mathbb C$ splitting $G_N$. Write
\[
 M=\bigoplus_{\chi,n}(\chi\otimes E(n))^{m_{\chi,n}},
 \qquad
 M'=\bigoplus_{\chi,n}(\chi\otimes E(n))^{m'_{\chi,n}}.
\]
Then $M\simeq M'$ if and only if
\[
 \det(1-TF_p\mid M)=\det(1-TF_p\mid M')
 \quad\text{for every prime }p\leq B_Q,\ p\nmid Q.                 \tag{AQ.1}
\]
Here $E(n)$ has Frobenius eigenvalue $p^{-n}$. These finitely many
polynomials recover every multiplicity $m_{\chi,n}$. The same
isomorphism criterion holds with rational coefficients.
The characters are labelled by their Euler-factor values
$\chi(p)$; changing the Artin-reciprocity convention inverts
all character labels and leaves the criterion unchanged.
One may use a single auxiliary $\ell>B_Q$ for all the
displayed Frobenius factors, avoiding the prime $\ell$
among the tested primes.

\paragraph{Known input and exact locator.}
Lamzouri, Li, and Soundararajan, \emph{Conditional bounds for the least
quadratic non-residue and related problems}, arXiv:1309.3595v3,
Corollary 1.2, p.~3, proves under Dirichlet GRH that every reduced residue
class $a\bmod Q$, $Q>3$, contains a prime
$p_a\leq(\varphi(Q)\log Q)^2$. This result is cited, not reproved.

\begin{proof}
The forward direction is functoriality. Inflate all characters to $G_Q$.
For each $a\in G_Q$ choose such a prime $p_a$.
Every eigenvalue at $p_a$ is $\chi(a)p_a^{-n}$, of complex absolute
value $p_a^{-n}$. Since $p_a>1$, different integers $n$ give different
absolute values. Equality in (AQ.1) therefore gives equality,
for every $a$ and $n$, of the sums of eigenvalues of absolute
value $p_a^{-n}$. Multiplying this sum by $p_a^n$ gives
\[
 t_n(a)=\sum_\chi m_{\chi,n}\chi(a).
\]
Finite character orthogonality yields the reconstruction
\[
 m_{\chi,n}=\frac1{|G_Q|}\sum_{a\in G_Q}
             t_n(a)\chi(a)^{-1}.                                \tag{AQ.2}
\]
Indeed, translating $\sum_a\theta(a)$ by an element $b$ with
$\theta(b)\ne1$ makes that sum zero for every nontrivial character
$\theta$; it is $|G_Q|$ for the trivial character. Thus (AQ.2)
gives the same multiplicities for $M'$.
For rational coefficients, the semisimple commutative algebra
$\mathbb Q[G_Q]$ is a product of cyclotomic fields. The
multiplicity of each simple rational module is recovered from
the common multiplicity of its characters over a splitting field.
The rational representations, weight by weight, are therefore
isomorphic as well.
\end{proof}

This is a finite stopping criterion for recognition within the
abelian Artin--Tate part of the programme. It needs every residue
class, or equivalent full character information: values on an
arbitrary generating set do not determine a general character.
It does not reconstruct extension classes, integral lattices,
period comparisons, or nonabelian fundamental groups.

\subsection{AQ-14-FORWARD: the cohomological implication at cote 14, scan 38}
\paragraph{Archive question.}
The sheet labels \emph{Th\'eor\`eme 3.1}, marks its proof
\emph{d\'em. pas faite}, and asks whether $u_Z=0$ implies,
modulo torsion, the vanishing of
\[
 a({}^tZ):H^{2(d-i)+1}(X,\mathbb Q_\ell(d-i))
                     \longrightarrow H^1(A,\mathbb Q_\ell).
                                                               \tag{AQ.3}
\]
Here $d=\dim X$ and $A$ is an abelian parameter variety.
Its margin asks \emph{A-t-on une r\'eciproque}.
The lower display $H^{(n-i)+1}$ is inconsistent with the
degree in the preceding paragraph; (AQ.3) uses the latter.

\paragraph{Precise modern formulation.}
For a smooth projective geometrically connected $X$ over an
algebraically closed field $k$, let $N^i(X)\subset CH^i(X)_{\rm alg}$
be generated by
\[
 \Gamma_*z,\quad
 z\in\ker\bigl(CH_0(T)^0\longrightarrow\operatorname{Alb}_T(k)\bigr),
 \quad\Gamma\in CH^i(T\times X),                                \tag{AQ.4}
\]
where $T$ is smooth projective geometrically connected.
Set $P^i(X)=CH^i(X)_{\rm alg}/N^i(X)$, the generated Albanese
quotient. For $Z\in CH^i(A\times X)$
the map on degree-zero cycles descends to
$u_Z:A(k)\to P^i(X)$, since two lifts of the same point of
$\operatorname{Alb}_A=A$ differ by an Albanese-kernel class.
If $u_Z=0$ integrally, then (AQ.3) is zero for every
$\ell\ne\operatorname{char}k$.
The sole initially supplied sheet 14/038 does not define
$A^i_{\rm alb}(X)$, but the official full cote now resolves
the definition. In its physical PDF page 24 (footer and
manuscript page 23), Definition 1.1 declares a cycle
Albanese-equivalent to zero when it has the form
$Z(v)$ for a smooth connected projective parameter $T$,
a degree-zero zero-cycle $v$ with $S_T(v)=0$, and a
correspondence $Z$ for which the intersection is defined.
The word \emph{projectif} is bracketed in the sheet,
between \emph{un sch\'ema} and \emph{lisse connexe}.
We first retain that projective-parameter formulation.
AQ-14-QUASIPROJECTIVE below proves equality with the
quasiprojective-parameter relation in characteristic zero,
with the nonproper Albanese and zero-cycle conventions explicit.
Here $S_T$ is the Albanese summation map.
Proposition 1.2(i), on the same page, says these cycles
form a subgroup containing rational equivalences and
defines the quotient of the Chow group by its image.
Proposition 1.2(v), physical PDF page 25 (footer 24),
says Albanese equivalence implies algebraic equivalence.
Finally Corollary 2.3, physical PDF page 34 (footer 33),
defines $P^i(X)$ itself as the algebraically trivial
subgroup of $A^i_{\rm alb}(X)$. Thus this is the
original quotient in the projective formulation,
not a quotient defined by regular maps to abelian
varieties. The Chow moving lemma makes the finitely
many intersections in (AQ.4) defined; since rational
equivalences are already killed, the cycle-level and
Chow formulations agree. Theorem 2.1, physical PDF
page 31 (footer 30), defines $u_Z$ through the
Albanese summation map. Its image lies in $P^i(X)$,
so vanishing in this subgroup and in the original
total quotient are equivalent. The supplied
\texttt{scan038} uses physical PDF numbering:
it is footer/manuscript page 37, after the cover.
Definition 1.1 itself says only that $k$ is a field.
The preceding handwritten sheets inspected in the full
cote concern other motive conventions and preparatory
cycle calculations; no explicit algebraic-closure
convention was found there. Accordingly the
algebraically closed hypothesis in our theorem is
retained as an explicit mathematical scope condition,
rather than attributed to an unseen original convention.

\paragraph{Primary inputs.}
Bloch (1979), \emph{Torsion algebraic cycles and a theorem of Roitman},
Propositions 3.5 (correspondences), 3.8 (specialization), and
Theorem 4.2 (prime-to-characteristic Roitman), provide the inputs.
They are also explicitly recorded in Achter, Casalaina-Martin,
and Vial, \emph{On the image of the second $\ell$-adic Bloch map},
arXiv:2007.07180v2, Propositions A.23, A.24, A.26, and A.29.
These appendix statements apply in every codimension.

\paragraph{Torsion lemma.}
For Bloch's map
\[
 \lambda_X^i:CH^i(X)[\ell^\infty]\longrightarrow
                    H^{2i-1}(X,\mathbb Q_\ell/\mathbb Z_\ell(i)),
\]
one has
\[
 N^i(X)\cap CH^i(X)[\ell^\infty]\subseteq\ker\lambda_X^i.          \tag{AQ.5}
\]
\begin{proof}
Choose $x$ in the intersection, a finite expression
$x=\sum_j\Gamma_{j*}z_j$ of the form (AQ.4), and $r$ with
$\ell^r x=0$. Spread out this entire expression, its
rational-equivalence certificates, its varieties and
correspondences, their Albanese maps, and the equalities
$\operatorname{alb}(z_j)=0$, over an integral finitely generated
ring over the prime ring. Shrink so that the varieties are
smooth projective, the chosen generic Albanese varieties extend
to abelian schemes $B_j$, and the chosen maps $T_j\to B_j$
extend. It is unnecessary to assert that these are universal
relative Albanese maps. Here the base can be chosen regular:
the finitely generated field of definition has a separating
transcendence basis over its prime field, and a model becomes
smooth over $\mathbb F_p$, or over $\mathbb Z[1/M]$ in
characteristic zero, after a further shrinking. All the
finitely many conditions just listed survive that shrinking.
In characteristic zero invert $\ell$ as well.
Choose a closed point with finite residue field of
characteristic $p\ne\ell$. A regular system of parameters
in its local ring gives a flag of prime ideals with regular
successive quotients. Localizing at consecutive primes
gives discrete valuation rings. Specialization along this
flag, after choosing algebraic closures of the residue
fields, gives geometric special fibres over
$\overline{\mathbb F}_p$, denoted by subscripts $0$.

Every degree-zero zero-cycle on a smooth projective connected
variety over $\overline{\mathbb F}_p$ is torsion: join its
finite support by finitely many projective curves and lift
the differences of points to their normalizations. Every
degree-zero divisor class on those curves is a point of a
Jacobian defined over a finite extension of a finite field,
and has finite order.

The specialized map $T_{j,0}\to B_{j,0}$ factors through
$h_j:\operatorname{Alb}(T_{j,0})\to B_{j,0}$. This map is an
isomorphism on $\ell$-power torsion. Indeed Kummer theory gives
$H^1(T,\mathbb Z_\ell(1))=T_\ell\operatorname{Pic}(T)
=T_\ell\operatorname{Pic}^0(T)_{\rm red}$: the N\'eron--Severi
group is finitely generated and has zero Tate module.
Duality of $\operatorname{Pic}^0_{\rm red}$ and Albanese
therefore gives
$H^1(T,\mathbb Z_\ell)=\operatorname{Hom}(T_\ell
\operatorname{Alb}_T,\mathbb Z_\ell)$, compatibly with maps.
This is also the codimension-one calculation in
Achter--Casalaina-Martin--Vial, Propositions A.25 and A.28.
The generic Albanese map thus induces an isomorphism on
$H^1(-,\mathbb Z_\ell)$, and smooth proper base change
preserves this isomorphism. The latter is obtained
by taking inverse limits of the finite-coefficient base-change
isomorphisms for each variety; no assertion that the Albanese
map is an isomorphism on every finite-coefficient $H^1$ is
needed. Equivalently $T_\ell h_j$ is an isomorphism.
It follows that $h_j$ is an isogeny with no $\ell$-primary
kernel, and hence an isomorphism on $\ell$-power torsion.
Since the specialization of the equality
$\operatorname{alb}(z_j)=0$ says that $z_{j,0}$ maps to zero
in $B_{j,0}$, the $\ell$-primary part of its actual Albanese
class is zero. Prime-to-$p$ Roitman implies that the
$\ell$-primary part of $z_{j,0}$ itself is zero. Thus every
$z_{j,0}$, and hence $x_0$, has order prime to $\ell$.
The equality $\ell^r x_0=0$ also specializes, so $x_0=0$.

Bloch's compatibility with specialization sends
$\lambda_X^i(x)$ to $\lambda_{X_0}^i(x_0)=0$. The
cohomological specialization map is an isomorphism by
smooth proper base change, first for finite coefficients
and then for their direct limit $\mathbb Q_\ell/\mathbb Z_\ell(i)$.
Consequently $\lambda_X^i(x)=0$. Only finitely many data are
spread for each $x$; no simultaneous model for all of $N^i(X)$
has been assumed.
\end{proof}

\begin{proof}[Proof of AQ-14-FORWARD]
Put $g=\dim A$. Roitman identifies
$CH_0(A)[\ell^\infty]$ with $A[\ell^\infty](k)$.
Use the torsion zero-cycle lift of a torsion point under this
isomorphism; it is unnecessary to identify it with the
literal class $[a]-[0]$. For every such torsion cycle $z$,
the premise gives $Z_*z\in N^i(X)\cap CH^i(X)[\ell^\infty]$.
Thus (AQ.5) gives $\lambda_X^i(Z_*z)=0$.
Correspondence compatibility, followed by taking Tate modules
and tensoring with $\mathbb Q_\ell$, makes
\[
 Z_*:H^{2g-1}(A,\mathbb Q_\ell(g))
             \longrightarrow H^{2i-1}(X,\mathbb Q_\ell(i))
\]
zero: the degree-zero Bloch map identifies the source
with $V_\ell A$. Under the perfect Poincar\'e pairings its
transpose is (AQ.3), so (AQ.3) is zero.
\end{proof}

This proves the stated generated-Albanese-quotient formulation
in every characteristic, independently of RH.

\paragraph{AQ-14-RATIONAL-PREMISE.}
The integral premise cannot be replaced by $u_Z\otimes\mathbb Q=0$.
Take an elliptic curve $E/\overline{\mathbb F}_p$ and
$X=A=E$, $Z=\Delta_E-E\times\{0\}$.
Then $P^1(E)=\operatorname{Pic}^0(E)$ and
$u_Z(a)=[a]-[0]$. The group $E(\overline{\mathbb F}_p)$ is
torsion, so the rationalized premise holds. The normalized
diagonal acts as the identity on the nonzero two-dimensional
$H^1(E,\mathbb Q_\ell)$, refuting the weakened assertion.
The phrase \emph{modulo torsion} in the original sheet
qualifies its cohomological conclusion.
For the asserted $P^1$ identification, a divisor
correspondence on $T\times E$ defines a morphism
$T\to\operatorname{Pic}^0(E)$ after choosing a base point,
and hence factors through $\operatorname{Alb}_T$.
Every generator of $N^1(E)$ is therefore zero.

\paragraph{AQ-14-FIELD-SCOPE.}
The algebraically closed base is also substantive.
If the premise is tested only on $k$-rational points
over an arbitrary field, it can fail to detect the
cohomological map. Take $k=\mathbb F_5$ and
$E:y^2=x^3+1$. The affine fibres over
$x=0,1,2,3,4$ have respectively $2,0,2,0,1$
points, so $\#E(\mathbb F_5)=6$, including infinity.
For $Z=6(\Delta_E-E\times\{0\})$, the map
$u_Z$ on $E(\mathbb F_5)$ is zero by Lagrange's
theorem, while the induced rational $\ell$-adic
map is $6\operatorname{id}\ne0$ for every $\ell\ne5$.
The original context's use of divisibility of
$A_Y(k)$ in physical PDF page 34 (footer 33)
is consistent with the algebraically closed
convention used here. This counterexample prevents
an unintended arbitrary-field reading.

\subsection{AQ-14-QUASIPROJECTIVE: removing the parameter qualification in characteristic zero}
Let $k$ be algebraically closed and let $X/k$ be smooth projective.
For a smooth connected parameter variety $T$, write
$Z_0(T)^0$ for the group of finite formal sums
$v=\sum_jn_j[t_j]$ of $k$-points with $\sum_jn_j=0$.
In this paragraph the Albanese of a nonproper $T$ means its
\emph{abelian Albanese}: the universal pointed morphism
$a_T:T\to\operatorname{Alb}_{\mathrm{ab}}(T)$ to an abelian
variety. Put $a_T(v)=\sum_jn_ja_T(t_j)$; this does not depend
on the base point used for the pointed morphism.

For $\Gamma\in\mathrm{CH}^i(X\times T)$ define
\[
 \Gamma(v)=\sum_j n_j\,\iota_{t_j}^{!}\Gamma
                         \in\mathrm{CH}^i(X),
\]
where $\iota_t:X\times\{t\}\hookrightarrow X\times T$
is the regular immersion of the fibre. Thus this is a
sum of fibre restrictions; no nonproper pushforward
$X\times T\to X$ is being used. Let $N^i_{\mathrm{qp}}(X)$
be generated by all such $\Gamma(v)$ with $T$ smooth connected
quasiprojective and $a_T(v)=0$. Let $N^i_{\mathrm{proj}}(X)$
use smooth connected projective $T$ instead.

If every permitted $T$ has a smooth projective compactification,
then
\[
 N^i_{\mathrm{qp}}(X)=N^i_{\mathrm{proj}}(X).                 \tag{AQ.14.1}
\]
In particular this equality holds in characteristic zero.
It also holds if the convention for a nonproper parameter is
the universal semiabelian Albanese, with the kernel imposed
on $Z_0(T)^0$. This last convention must be formulated on
zero-cycles, not automatically on the ordinary Chow group
of the nonproper variety.

\begin{proof}
The inclusion from right to left is immediate, since a
smooth projective parameter is quasiprojective. For the
other inclusion take a generator $\Gamma(v)$ as above.
Choose a smooth projective compactification
$j:T\hookrightarrow\overline T$. Represent $\Gamma$ by
a finite integral linear combination of integral
codimension-$i$ subvarieties of $X\times T$, and take
their closures in $X\times\overline T$. This gives a class
$\overline\Gamma\in\mathrm{CH}^i(X\times\overline T)$
whose restriction is $\Gamma$.

The morphism $a_{\overline T}\circ j:T\to
\operatorname{Alb}(\overline T)$, normalized to send the
chosen base point to zero, factors by universality through
$a_T$. Therefore
$a_{\overline T}(j_*v)=0$. Here $j_*v$ denotes the same
finite formal zero-cycle regarded on $\overline T$;
it does not assert the existence of a pushforward
$j_*:\mathrm{CH}_0(T)\to\mathrm{CH}_0(\overline T)$.
For each point in the support of $v$, refined fibre
restriction commutes with restriction to the open subset
$X\times T$; equivalently, both restrictions are computed
in the same neighbourhood of $X\times\{t_j\}$. Hence
\[
 \overline\Gamma(j_*v)=\Gamma(v).
\]
The left side is a generator allowed in
$N^i_{\mathrm{proj}}(X)$, proving the reverse inclusion.
The same proof works for the semiabelian convention,
because the target $\operatorname{Alb}(\overline T)$ is
also semiabelian, and because the semiabelian Albanese
of a smooth projective parameter is its ordinary
abelian Albanese.

In characteristic zero, embed the quasiprojective $T$
in a projective variety and resolve that projective
closure by a projective resolution that is an
isomorphism over its smooth locus, in particular over
$T$. The resulting $\overline T$ is smooth projective.
This is the strong compactification consequence of
Hironaka's resolution theorem, stated in precisely
this form in Kahn and Sujatha, \emph{Birational geometry
and localisation of categories}, May 22, 2008 author
version, Theorem 1.6.3(a), p.~11.
\end{proof}

Thus AQ-14-FORWARD applies in characteristic zero to the
quotient generated using all smooth quasiprojective
parameters, and AQ-14-CONVERSE over $\mathbb C$ remains
a counterexample for that larger parameter class as
well. The algebraically closed base-field qualification
remains. In positive characteristic this comparison
applies to parameters admitting a smooth projective
compactification; no existence of such compactifications
in all dimensions is asserted here. The next theorem proves the
rational cohomological forward implication without needing that
positive-characteristic equality of relations. AQ-14-ALTERATION then proves that the discrepancy between those relations is $p$-primary torsion.

\subsection{AQ-14-NONPROPER: the rational cohomological implication for all smooth quasiprojective parameters}
Retain the algebraically closed base $k$ and smooth projective $X$.
Let $N^i_{\mathrm{qp}}(X)$ be the relation defined in
AQ-14-QUASIPROJECTIVE, using the \emph{abelian} Albanese of a
smooth connected quasiprojective parameter and free degree-zero
zero-cycles. Put
$P^i_{\mathrm{qp}}(X)=\mathrm{CH}^i(X)_{\mathrm{alg}}/
N^i_{\mathrm{qp}}(X)$. For every prime
$\ell\ne\operatorname{char}k$, one has
\[
 N^i_{\mathrm{qp}}(X)\cap\mathrm{CH}^i(X)[\ell^\infty]
       \subseteq\ker\lambda_X^i.                            \tag{AQ.14.2}
\]
Consequently the implication of original 14/038 holds for this
quasiprojective-parameter quotient in every characteristic:
if $A$ is abelian, $Z\in\mathrm{CH}^i(A\times X)$, and
$u_Z:A(k)\to P^i_{\mathrm{qp}}(X)$ is zero integrally, then
\[
 a({}^tZ):H^{2(\dim X-i)+1}(X,\mathbb Q_\ell(\dim X-i))
                       \longrightarrow H^1(A,\mathbb Q_\ell)
\]
is zero. The same implication holds for the smaller relation
using the semiabelian Albanese kernel. No equality
$N^i_{\mathrm{qp}}=N^i_{\mathrm{proj}}$ in arbitrary positive
characteristic is asserted.

\paragraph{Inputs and coefficient conventions.}
Write $H_j^S(T,R)$ for Suslin homology: it is the homology of
the complex generated in degree $n$ by integral closed
subvarieties of $T\times\Delta^n$ finite and surjective over
$\Delta^n$, tensored with $R$. Its degree-zero cycles are
finite formal sums of closed points. Write $G_T$ for the
pointed Serre semiabelian Albanese and
\[
 0\longrightarrow S_T\longrightarrow G_T\longrightarrow B_T
       \longrightarrow0
\]
for its torus and abelian quotient. Universality shows that
$B_T$ is the abelian Albanese of $T$: a homomorphism from a
torus to an abelian variety is zero. The notation $[\ell^\infty]$
means the subgroup annihilated by some power of $\ell$.

Geisser, \emph{Rojtman's theorem for normal schemes},
arXiv:1402.1831v2, Theorem 1.1, p.~2, gives, for normal $T$,
\[
 H_0^S(T,\mathbb Z)[\ell^\infty]
       \xrightarrow{\sim}G_T(k)[\ell^\infty],\qquad
 H_1^S(T,\mathbb Z)\otimes\mathbb Q_\ell/\mathbb Z_\ell=0.
                                                               \tag{AQ.14.3}
\]
The degree-zero convention on the left is automatic for
torsion. This prime-to-characteristic theorem does not
require a smooth compactification. The finite-coefficient
perfect pairing used below is Geisser--Schmidt (2016),
Theorem 1.1, journal p.~92; their Theorem 7.1, p.~118, identifies it
with the Suslin--Voevodsky comparison. We invoke its
prime-to-characteristic case, where tame and ordinary
abelian $\ell$-primary covers agree. Kahn (2014),
\emph{Algebraic tori as Nisnevich sheaves with transfers},
Lemma 4.1(b), journal p.~707 (July 4, 2013 author version,
Lemma 3(b), p.~8), states
$\operatorname{Alb}_{\mathrm{sa}}(G)=G$ for a semiabelian $G$.

The integral generic-base-change input is Sun,
arXiv:1511.01393v1, Definition 2.2, p.~3, Lemma 2.3(iii),
pp.~4--5, and Theorem 2.5(i), p.~7, with
$\Lambda=\mathbb Z_\ell$. In Lemma 2.3(iii)'s proof, one
dense open works for \emph{all} finite coefficient levels,
by the finite set of Jordan--H\"older constituents. The
constant $\mathbb Z_\ell$ system has the single simple
constituent $\mathbb F_\ell$. Constructibility and preservation
of stratifiability are part of the stated setup and
Remark 2.5.1(ii). We use this only over a finite-type
$\mathbb F_p$-base, $p\ne\ell$, which satisfies the finite
cohomological-dimension condition (LO).

\begin{proof}
Characteristic zero follows already from
AQ-14-QUASIPROJECTIVE and AQ-14-FORWARD. Thus assume
$\operatorname{char}k=p>0$ and $\ell\ne p$.
We first give the three facts needed for specialization.

\emph{1. A correspondence with proper target descends to
Suslin zero-homology.}
For $\Gamma\in\mathrm{CH}^i(X\times T)$, the fibre-evaluation
map $Z_0(T)\to\mathrm{CH}^i(X)$ kills every Suslin boundary.
Indeed, a generator of a Suslin one-chain is an integral
$W\subset\mathbb A^1\times T$ finite and surjective over
$\mathbb A^1$. Its normalization $C$ is a smooth curve,
finite over $W$ and over $\mathbb A^1$. Let $\overline C$
be its smooth projective completion. Pull $\Gamma$ back
along the morphism $X\times C\to X\times T$; pullback is
well defined between smooth varieties by the graph/refined
Gysin construction. Extend a representing cycle by closure
to a class on $X\times\overline C$.

The function $t/(t-1)$ on $\overline C$ has divisor equal
to the fibre over zero minus the fibre over one, including
multiplicities, and no contribution over infinity. The
normalization projection formula identifies their images
in $Z_0(T)$ with the two faces of $W$. Applying the extended
correspondence to this principal divisor gives zero in
$\mathrm{CH}^i(X)$. Its evaluations at points of $C$ agree
with the original fibre evaluations, since refined pullback
is local around those fibres. This proves the claim and
produces a homomorphism
\[
 \Gamma_*:H_0^S(T,\mathbb Z)\longrightarrow\mathrm{CH}^i(X).
                                                               \tag{AQ.14.4}
\]
The construction is compatible with extension of the base
field and with automorphisms of that field. No pushforward
along the nonproper projection $X\times T\to X$ was used.

\emph{2. The integral $H^1$ of the semiabelian Albanese.}
The homology coefficient sequence and (AQ.14.3) give a
natural isomorphism
\[
 H_1^S(T,\mathbb Q_\ell/\mathbb Z_\ell)
             \simeq G_T(k)[\ell^\infty].                    \tag{AQ.14.5}
\]
For clarity, the coefficient sequence has left term
$H_1^S(T,\mathbb Z)\otimes\mathbb Q_\ell/\mathbb Z_\ell$
and right term
$\operatorname{Tor}_1^{\mathbb Z}(H_0^S(T,\mathbb Z),
\mathbb Q_\ell/\mathbb Z_\ell)=H_0^S(T,\mathbb Z)[\ell^\infty]$;
the left term is zero by (AQ.14.3).

The prime-to-$p$ finite pairings of Geisser--Schmidt are
natural and compatible with the coefficient homomorphisms.
Dualizing the direct system
$H_1^S(T,\mathbb Z/\ell^n)$, whose transition is induced by
$\mathbb Z/\ell^n\hookrightarrow\mathbb Z/\ell^{n+1}$,
identifies its dual with the inverse system
$H^1_{\mathrm{et}}(T,\mathbb Z/\ell^n)$ under reduction.
The pairings are viewed in $\mathbb Q_\ell/\mathbb Z_\ell$
via $1\bmod\ell^n\mapsto\ell^{-n}\bmod\mathbb Z_\ell$;
this makes the transition compatibility explicit. Thus
\[
 H^1_{\mathrm{et}}(T,\mathbb Z_\ell)
  \simeq\operatorname{Hom}\bigl(
       G_T(k)[\ell^\infty],\mathbb Q_\ell/\mathbb Z_\ell\bigr)
  \simeq\operatorname{Hom}_{\mathbb Z_\ell}
       (T_\ell G_T,\mathbb Z_\ell).                         \tag{AQ.14.6}
\]
The left side is the inverse limit of the finite groups;
there is no $\varprojlim^1 H^0$ obstruction, because the
constant $H^0$ transition maps are surjective. All
isomorphisms are functorial. Since $G_T$ is its own Serre
Albanese, the Albanese morphism $a_T:T\to G_T$ induces an
isomorphism on $H^1(-,\mathbb Z_\ell)$. This conclusion
is deliberately integral-adic, not an assertion about
$H^1(-,\mathbb Z/\ell^n)$ for each $n$.

\emph{3. The finite-field vanishing.}
Suppose for the moment that $k=\overline{\mathbb F}_p$.
Every class in $H_0^S(T,\mathbb Z)^0$ is torsion.
To see this, join the finitely many points in a representative
by finitely many integral curves in the connected smooth
quasiprojective variety $T$, and pass to their normalizations.
Such curves can be chosen through any prescribed finite
set by hyperplane sections of sufficiently high degree;
only the resulting curves and their normalizations are
used. Differences of points lift to degree-zero classes
on these smooth open curves. For a smooth curve $C$,
$H_0^S(C,\mathbb Z)^0$ is the group of points of its
Rosenlicht generalized Jacobian, equivalently the curve
case of Geisser, Theorem 1.3, p.~3. Each such Jacobian and
each of its points are defined over some finite extension
of a finite field, and its rational points over that
finite field form a finite group. Each lifted class,
and hence each original degree-zero class, is torsion.

Let $v\in H_0^S(T,\mathbb Z)^0$ have zero image in $B_T(k)$,
and let $v_\ell$ denote its $\ell$-primary component.
Under (AQ.14.3), the class $v_\ell$ corresponds to an
element of $S_T(k)[\ell^\infty]$. Consider the homomorphism
\[
 S_T(k)[\ell^\infty]\longrightarrow
 H_0^S(T,\mathbb Z)[\ell^\infty]
 \xrightarrow{\Gamma_*}\mathrm{CH}^i(X)[\ell^\infty]
 \xrightarrow{\lambda_X^i}
 H^{2i-1}(X,\mathbb Q_\ell/\mathbb Z_\ell(i)).                \tag{AQ.14.7}
\]
All varieties, maps and $\Gamma$ descend to a common finite
field $\mathbb F_q$, after enlarging it if needed. Thus
(AQ.14.7) is Frobenius-equivariant. On
$V_\ell S_T=T_\ell S_T\otimes\mathbb Q_\ell$ geometric
Frobenius has eigenvalues $q^{-1}$ times roots of unity:
a torus splits over a finite extension of $\mathbb F_q$.
On
$H^{2i-1}(X,\mathbb Q_\ell(i))$ its eigenvalues have
complex absolute value $q^{-1/2}$, by Deligne (1974),
Theorem (1.6), printed p.~276, after the indicated Tate
twist. These spectra are disjoint. A linear map
commuting with Frobenius between these two spaces is zero:
apply relatively prime characteristic polynomials and
B\'ezout's identity, or use generalized eigenspaces
after extending scalars. Semisimplicity is not
required for this step.

The rational Tate module of the last group in (AQ.14.7)
is exactly $H^{2i-1}(X,\mathbb Q_\ell(i))$. More explicitly,
that group has maximal divisible subgroup
$H^{2i-1}(X,\mathbb Z_\ell(i))_{\mathrm{free}}
 \otimes\mathbb Q_\ell/\mathbb Z_\ell$ and finite quotient
contained in $H^{2i}(X,\mathbb Z_\ell(i))[\ell^\infty]$.
The domain of (AQ.14.7) is divisible. Its image is therefore
in the maximal divisible subgroup, and a homomorphism
between these finite-corank divisible $\ell$-groups is
zero if its map of rational Tate modules is zero.
Consequently (AQ.14.7) is zero, and
$\lambda_X^i(\Gamma_*v_\ell)=0$.

It follows that over $\overline{\mathbb F}_p$ the conclusion
(AQ.14.2) holds: in a finite sum of generators each
$H_0^S$ class is torsion; its prime-to-$\ell$ component
maps to prime-to-$\ell$ torsion, whereas the $\ell$-primary
component has zero Bloch image by the preceding argument.
Taking the $\ell$-primary component of an expression for
$x\in N^i_{\mathrm{qp}}(X)\cap\mathrm{CH}^i(X)[\ell^\infty]$
gives the claimed vanishing.

\emph{4. Spreading to the finite-field situation.}
Now take arbitrary algebraically closed $k$ of characteristic
$p>0$ and a fixed
$x\in N^i_{\mathrm{qp}}(X)\cap\mathrm{CH}^i(X)[\ell^\infty]$.
Choose a finite expression
$x=\sum_j\Gamma_{j*}v_j$, with $v_j\in Z_0(T_j)^0$
mapping to zero in $B_{T_j}(k)$, and choose $r$ with
$\ell^r x=0$. Spread all varieties, finite zero-cycles,
correspondences, rational-equivalence witnesses for these
equalities, and the pointed maps
$T_j\to G_{T_j}\to B_{T_j}$ over an integral finite-type
$\mathbb F_p$-scheme $S$. After shrinking, $S$ is regular,
$X/S$ is smooth projective, the $T_j/S$ are smooth with
geometrically connected fibres, and the $G_j/S$ are
semiabelian schemes in exact sequences
$0\to S_j\to G_j\to B_j\to0$ with tori $S_j$ and
abelian schemes $B_j$. The sections representing the
zero-cycles remain inside $T_j$, and the relations
$B_j(v_j)=0$ persist. The spread maps need not be
relative universal Albanese morphisms.

Apply Sun's integral generic base change simultaneously to
the finitely many $T_j\to S$ and $G_j\to S$. Their
$R^1(-)_*\mathbb Z_\ell$ are constructible adic sheaves,
and become lisse on a further dense open. By (AQ.14.6)
the map induced by $T_j\to G_j$ is an isomorphism at
a geometric generic point. Indeed its base extension
to $k$ is the chosen Albanese map, and prime-to-$p$
\'etale cohomology is unchanged by extension of
algebraically closed fields. Its kernel and cokernel
vanish on a further dense open, so this is an
isomorphism of integral adic sheaves there. Generic
base change now identifies its stalk at \emph{every}
geometric point of that open with the actual map
on $H^1(-,\mathbb Z_\ell)$ of the fibre. This uses one
open for the full integral system, not separate opens
for infinitely many coefficient levels.

Choose a closed point $s$ in this open; its residue
field is finite. Over $\overline{k(s)}$, the map
$T_{j,s}\to G_{j,s}$ factors by universality as
\[
 T_{j,s}\longrightarrow G_{T_{j,s}}
                  \xrightarrow{h_j}G_{j,s}.
\]
By (AQ.14.6), $T_\ell h_j$ is an isomorphism. A
homomorphism of semiabelian varieties whose Tate module
is an isomorphism is an isogeny with no $\ell$-primary
kernel: a positive-dimensional connected kernel or
quotient has a nonzero Tate module, and a nonzero finite
$\ell$-primary kernel makes the Tate-module image a
proper sublattice. The induced homomorphism of abelian
quotients
\[
 b_j:B_{T_{j,s}}\longrightarrow B_{j,s}
\]
is also an isogeny with no $\ell$-primary kernel.
Indeed an isogeny of semiabelian varieties maps the
torus onto the torus; the kernel of its map of abelian
quotients is a quotient of its finite kernel. The first
assertion follows since the resulting quotient of the
target by the image torus is a quotient of the proper
source abelian variety, and hence has zero toric part.

The equality $B_{j,s}(v_{j,s})=0$ and injectivity of
$b_j$ on $\ell$-primary torsion imply that the
$\ell$-primary component of the actual abelian Albanese
class of $v_{j,s}$ is zero. Although the whole class
need not vanish, this is precisely what step 3 needs:
the $\ell$-primary Suslin class of $v_{j,s}$ corresponds
to the actual torus $S_{T_{j,s}}$. Thus the argument
in step 3 gives $\lambda_{X_s}^i(x_s)=0$.

Finally specialize from a geometric generic point to
$s$ along a regular flag through $s$, as in the complete
flag construction in AQ-14-FORWARD. The finite cycle
and rational-equivalence data make $x_s$ the
specialization of $x$, preserve $\ell^r x_s=0$, and
identify each specialized fibre evaluation with
$\Gamma_{j,s*}v_{j,s}$. Refined Gysin functoriality
provides this compatibility locally at the sections.
Bloch's specialization theorem sends
$\lambda_X^i(x)$ to $\lambda_{X_s}^i(x_s)$, and smooth
proper base change for $X/S$ makes the cohomological
specialization an isomorphism. Cohomological invariance
under algebraically closed extension also identifies
the original $k$-fibre with the geometric generic
fibre for this calculation. Therefore
$\lambda_X^i(x)=0$, proving (AQ.14.2).

For the final cohomological implication, lift each
$\ell$-power torsion point of $A$ to its torsion
zero-cycle class by prime-to-$p$ Roitman. The integral
premise puts its image under $Z$ in the left side of
(AQ.14.2). Bloch correspondence compatibility, followed
by Tate modules and tensoring with $\mathbb Q_\ell$,
therefore kills
$Z_*:H^{2\dim A-1}(A,\mathbb Q_\ell(\dim A))
\to H^{2i-1}(X,\mathbb Q_\ell(i))$.
Poincar\'e duality gives exactly the displayed transpose.
This is the same final argument as AQ-14-FORWARD, now
with the broader relation. For the semiabelian-kernel
convention, its relation is contained in the
abelian-kernel relation, so the torsion vanishing and
the implication follow as well.
\end{proof}

\subsection{AQ-14-ALTERATION: the two parameter relations agree away from the characteristic}
\paragraph{Statement.}
Let $k$ be algebraically closed of characteristic $p>0$,
and let $X/k$ be smooth projective.  Define
$N^i_{\rm proj}(X)\subset N^i_{\rm qp}(X)\subset CH^i(X)$
using the projective and smooth quasiprojective parameters,
respectively, and the ordinary abelian Albanese, exactly as
in AQ-14-QUASIPROJECTIVE.  Then
\[
 N^i_{\rm qp}(X)/N^i_{\rm proj}(X)
 \quad\hbox{is a $p$-primary torsion group.}               \tag{AQ.ALT.1}
\]
In particular the relations agree after tensoring with
$\mathbb Z[1/p]$, hence also with $\mathbb Q$.
The natural map between the two cycle quotients has
$p$-primary kernel and induces an isomorphism on
$\ell$-primary torsion for every $\ell\ne p$.
The same assertions hold if the nonproper relation uses
zero cycles annihilated by the semiabelian Albanese.
Together with AQ-14-QUASIPROJECTIVE this gives equality
with rational coefficients in every characteristic.
No bounded exponent for the possible $p$-primary kernel
is asserted.

\paragraph{Exact external inputs.}
The separability refinement is needed here: Illusie and
Temkin (2014), Expos\'e X, Theorem 3.5(iii), printed p.~202,
gives a \emph{separable} projective prime-to-$\ell$ alteration
with regular source over a perfect field.  Their
Theorem 2.1 alone is not the separability statement;
the distinction is explicit on p.~167 and in \S3.5.4,
p.~207.  The other input is the extension of a rational
map from a smooth variety to an abelian variety to a
morphism on the whole source: Milne, \emph{Abelian
Varieties}, version 2.0 (March 16, 2008), Chapter I,
Theorem 3.2, printed p.~17.  We use these as cited
classical inputs.  The moving and trace arguments
needed for their application are proved below.

\begin{proof}
\emph{1. Moving zero cycles inside a dense open using Suslin relations.}
For a smooth connected quasiprojective $T$ and a nonempty
open $U\subset T$, every class in $H^S_0(T,\mathbb Z)$
is represented by a zero cycle supported in $U$.
The replacement can be chosen with the same degree.
We prove this directly on point generators.
There is nothing to prove for a point already in $U$,
or if $\dim T=0$.
Fix $x\in T\setminus U$.  An affine neighborhood of $x$
embeds into affine space.  Successive general hyperplanes
through $x$ cut its dimension down to one and cut
$T\setminus U$ down to dimension zero in that neighborhood.
At each step one avoids finitely many positive-dimensional
components; this is possible because $k$ is infinite and
the linear system has only the point $x$ as its base locus.
An irreducible one-dimensional component through $x$
therefore meets $U$.  Its closure in $T$ is an integral
curve $C_0$ through $x$, still meeting $U$.
Let $C\to C_0$ be its normalization and choose
$x'\in C$ above $x$.  The normalization is finite,
and $C$ is a smooth curve because $k$ is perfect.
Let $\overline C$ be its smooth projective completion,
$D=\overline C\setminus C$, and
$F=C\setminus(C\times_T U)$.
Both $D$ and $F$ are finite sets.  Choose
$y\in C\times_T U$.

For sufficiently large $n$, Riemann--Roch gives a rational
function $f$ on $\overline C$, with poles only at $y$,
whose first-order jet at $x'$ is a uniformizer and whose
value at every point of $D\cup(F\setminus\{x'\})$
is $1$.  Explicitly the required evaluation map is
surjective because the first cohomology of
\[
 \mathcal O_{\overline C}
 \left(ny-2x'-\sum_{z\in D\cup(F\setminus\{x'\})}z\right)
\]
vanishes once its degree exceeds $2g(\overline C)-2$.
Thus $f$ has a simple zero at $x'$ and no zero or pole
at the other points of $D\cup F$.
Its divisor is of the form
\[
 \operatorname{div}(f)=[x']+z-a[y],
 \qquad \operatorname{supp}(z)\subset C\times_T U,
\]
where $z$ is effective and $a>0$.
This divisor is a Suslin boundary on $C$:
the nonconstant morphism
$g=f/(f-1):\overline C\to\mathbb P^1$ is finite,
and $g^{-1}(\mathbb A^1)\subset C$ since $f=1$ on $D$.
Its graph gives a finite correspondence from
$\mathbb A^1$ to $C$; its fibres at $0$ and $1$
differ by $\operatorname{div}(f)$, up to the harmless
choice of boundary sign.  Push this correspondence to
$T$.  It shows that $[x]$ has a same-degree representative
supported in $U$.  Linearity proves the moving assertion.

Both maps used in our relation respect these Suslin
boundaries.  For the correspondence action
$\Gamma_*:Z_0(T)\to CH^i(X)$, this is the proper-target
Suslin descent proved in Step 1 of AQ-14-NONPROPER:
pull to a normalized curve, close on its projective
completion, and evaluate the principal divisor there.
For the ordinary Albanese map it follows directly as well:
the morphism $C\to\operatorname{Alb}_{\rm ab}(T)$
extends to $\overline C$ by properness, and any principal
divisor has zero sum in that abelian variety.
Consequently, if $v$ has degree zero and Albanese sum zero,
we may replace it by a zero cycle $v_U$ supported in $U$
with
\[
 \deg v_U=0,\qquad\operatorname{alb}_T(v_U)=0,
 \qquad \Gamma_*(v_U)=\Gamma_*(v).                \tag{AQ.ALT.2}
\]

\emph{2. The prime-to-$\ell$ transfer.}
Fix a prime $\ell\ne p$ and a generator
$\Gamma_*(v)$ of $N^i_{\rm qp}(X)$.
Choose a projective integral compactification
$T\subset\overline T$.
The cited Theorem 3.5(iii), applied over the perfect field
$k$, gives a regular projective source and a separable
projective prime-to-$\ell$ alteration of $\overline T$.
Choose a connected component that dominates
$\overline T$; write it as
\[
 f:S\longrightarrow\overline T,\qquad
 d=[k(S):k(T)],\qquad (d,\ell)=1.
\]
It is smooth over $k$, projective and integral.
There exists a nonempty open $U\subset T$ over which
$f^{-1}(U)\to U$ is finite \'etale of degree $d$.
Use Step 1 to obtain $v_U$ satisfying (AQ.ALT.2).

Fix a point of $S$ to normalize
$a_S:S\to A_S=\operatorname{Alb}(S)$.
The fibre sum
\[
 \nu:U\longrightarrow A_S,\qquad
 u\longmapsto\sum_{s\in f^{-1}(u)}a_S(s)
\]
is a morphism: on an \'etale cover of $U$ the finite
\'etale cover splits into $d$ sections, their sum is
regular, and these sums descend because they are
invariant under permutations of the sections.
It is a rational map from smooth $T$ to an abelian
variety, so the cited extension theorem gives a morphism
$\overline\nu:T\to A_S$.
After subtracting its value at a base point, universality
of $\operatorname{Alb}_{\rm ab}(T)$ makes this morphism
factor through a homomorphism
$\operatorname{Alb}_{\rm ab}(T)\to A_S$.
Write $v_U=\sum_u n_u[u]$.  Since it has degree and Albanese sum zero,
\[
 w=f_U^*v_U\in Z_0(S)^0,
 \qquad \operatorname{alb}_S(w)
        =\sum_u n_u\nu(u)=0.                       \tag{AQ.ALT.3}
\]

Restrict a cycle representing $\Gamma$ to $X\times U$,
pull it back along the finite \'etale map
$X\times f^{-1}(U)\to X\times U$, and take its
cycle-theoretic closure on $X\times S$.
Call the resulting codimension-$i$ cycle $\Gamma'$.
Fibre Gysin evaluation commutes with this \'etale
pullback and with restriction to the open set.
Each of the $d$ points above $u\in U(k)$ therefore
has the same fibre class $\Gamma_u\in CH^i(X)$.
It follows that
\[
 \Gamma'_*(w)=d\,\Gamma_*(v_U)
             =d\,\Gamma_*(v)
             \in N^i_{\rm proj}(X),               \tag{AQ.ALT.4}
\]
where membership uses the projective parameter $S$
and (AQ.ALT.3).

\emph{3. The quotient and its torsion.}
Every element $c$ of $N^i_{\rm qp}/N^i_{\rm proj}$
is a finite sum of such generators.  For each
$\ell\ne p$, take the product of their degrees in
(AQ.ALT.4); it is prime to $\ell$ and annihilates $c$.
Taking any one $\ell\ne p$ first shows that $c$ is
torsion, of some finite order $n$.  For every
$\ell\ne p$ a prime-to-$\ell$ annihilator forces
$\ell\nmid n$.  Thus $n$ is a power of $p$, proving
(AQ.ALT.1).

Put $Q_{\rm proj}=CH^i(X)/N^i_{\rm proj}$ and
$Q_{\rm qp}=CH^i(X)/N^i_{\rm qp}$, and let
$K=N^i_{\rm qp}/N^i_{\rm proj}$.
The exact sequence
$0\to K\to Q_{\rm proj}\to Q_{\rm qp}\to0$
has $p$-primary torsion kernel.  Multiplication by
$\ell$ is an automorphism of $K$ when $\ell\ne p$:
this holds on every finite cyclic subgroup and hence
on their union.  Therefore the induced map on
$\ell$-primary torsion is injective.  It is surjective
as well: if $y\in Q_{\rm qp}$ is killed by
$\ell^a$, lift it to $x\in Q_{\rm proj}$;
then $\ell^a x\in K$, and choose $z\in K$ with
$\ell^a z=\ell^a x$.  The lift $x-z$ is killed by
$\ell^a$.  This proves the assertion about torsion.
Localization is flat and kills $K$, proving the
claimed equalities after inverting $p$.

Finally, a zero cycle in the semiabelian Albanese
kernel is in the abelian Albanese kernel.  Its relation
is therefore between $N^i_{\rm proj}$ and the relation
just treated.  The same $p$-primary conclusion follows
for that intermediate subgroup and quotient.
\end{proof}

\paragraph{Consequence for the original cohomological question.}
This also gives a second route to the forward implication
of 14/038.  If an $\ell$-primary torsion cycle
$x\in CH^i(X)$ belongs to $N^i_{\rm qp}$, (AQ.ALT.4),
applied to a finite expression for $x$, supplies
$d$ prime to $\ell$ with $dx\in N^i_{\rm proj}$.
AQ-14-FORWARD gives $d\lambda_X^i(x)=0$.
Since $\lambda_X^i(x)$ is $\ell$-primary and
multiplication by $d$ is invertible on that group,
$\lambda_X^i(x)=0$.
The original integral-premise and rational-transpose
argument then applies.  This deduction uses the
projective torsion lemma and the alteration theorem;
the earlier direct Suslin-weight proof is retained
as a separately checked argument.

\subsection{AQ-14-CONVERSE: the generated-quotient converse fails}
For the generated quotient just defined, over $\mathbb C$
the converse fails, even rationally in its target.
The following is an application of Walker (2007), Theorem 5.7(4),
p.~930, and Theorems 6.1 and 6.2, pp.~933--934.
Take a smooth split six-dimensional quadric $W$,
a sufficiently general high-degree smooth hypersurface section
$Y\subset W$, and the two-cycle $\gamma$ supplied by Theorem 6.1
with $r=2$: its singular homology class is zero, whereas
$s\gamma\ne0$ in $L_1H_4(Y,\mathbb Q)$.
The two families of maximal $\mathbb P^3$'s on $W$
give independent middle-dimensional algebraic
homology classes, supplying the rank-at-least-two
hypothesis of Theorem 6.1.
Here $\dim Y=5$. For any complex elliptic curve $E$ put
$X=Y\times E$, so $\dim X=6$, and
\[
 Z=\gamma\times(\Delta_E-E\times\{0\})\in CH^4(E\times X).
\]
Then (AQ.3) is zero but $u_Z$ is nonzero in $P^4(X)$,
even after tensoring with $\mathbb Q$.
\begin{proof}
The cycle class of $Z$ is zero because that of $\gamma$ is zero,
so its cohomological correspondence vanishes.
The needed factorization through the generated Albanese
quotient follows explicitly from Walker's Theorem 5.7(2),(4),
p.~930, and Example 5.5, p.~928. If $d=\dim X$,
$r=d-i$, and $\Gamma\in CH^i(T\times X)$, those statements give
\[
 \Phi^{\rm mor}_{X,r}(\Gamma_*z)
    =\Gamma_{J,*}\Phi^{\rm mor}_{T,0}(z)
    =\Gamma_{J,*}\operatorname{alb}_T(z).
\]
Theorem 5.7(4) applies to every Chow correspondence, not just
an equidimensional representative. The displayed expression
vanishes for each generator (AQ.4). Thus
$\Phi^{\rm mor}_{X,r}$ descends to $P^i(X)$.
The product construction in the proof of Walker's Theorem 6.2
gives a nonzero cylinder
\[
 \gamma\otimes H_1(E,\mathbb Q)\longrightarrow L_2H_5(X,\mathbb Q).
\]
Pairing its image with a dual class on $E$ and pushing to $Y$
gives a nonzero multiple of $s\gamma$, verifying nonvanishing.
This is a morphism of the (ind-)Hodge structures defining
the morphic Jacobian. The source has type $(-1,0),(0,-1)$
after twisting. A nonzero rational morphism from this source
has nonzero tangent map on the associated Jacobians: if the
relevant complex Hodge summand were killed, its conjugate
would also be killed, so the whole morphism would vanish.
More explicitly, strictness of the Hodge filtration gives
a nonzero complex linear map after quotient by $F^0$;
if it were zero, the whole real image would lie in
$F^0\cap\overline{F^0}$, which is zero in weight $-1$.
The underlying Lawson group is countable, as used in
Walker's proof of Theorem 6.2, p.~934. Quotienting the
nonzero vector-space image by its intersection with
this countable period group still gives uncountably
many image points. Torsion in a complex vector space
modulo a countable subgroup is countable, because
its inverse image is the countable union of the
groups $\tfrac1n\Lambda$. Hence
$a\mapsto\gamma\times([a]-[0])$ has nontorsion image
in the morphic Jacobian.
Since it factors through $P^4(X)$, the converse fails
even rationally.
\end{proof}
No unverified effective degree, such as degree 14, is
needed: the primary theorem's sufficiently-large-degree
construction supplies the counterexample. This is a
translation of a known phenomenon to the original
projective quotient of 14/038, independent of RH.
Walker's morphic Jacobian
need not be a finite-dimensional abelian variety.
Consequently this argument does not refute the separate
universality question for the classical Abel--Jacobi
map among finite-dimensional regular abelian targets;
Walker discusses that distinction in \S7.

\subsection{AQ-WEIL-EXIST: the strong conjecture at 15/100}
The source says that its preceding necessary conditions
are sufficient, following the boxed condition
$|\sigma\alpha|=q^{i/2}$. In the unambiguous modern form,
every nonzero algebraic integer with these absolute values
occurs on $H^i$ of a smooth projective scheme over
$\mathbb F_q$. This is true for every $i\geq0$.

The unconditional existence argument is the surjectivity
part of Milne (1994), Proposition 2.6, p.~15 of the author's
manuscript. The whole proposition classifies simple motives
under the Tate conjecture, explicitly assumed at the start
of \S2, p.~12; its full classification must retain that
assumption. Honda (1968), \emph{Isogeny classes of abelian
varieties over finite fields}, Main Theorem, p.~83, is
the unconditional existence input.
\begin{proof}[Short proof of the specified $H^i$ form]
For $i>0$ choose $\beta$ with $\beta^i=\alpha$.
It is integral over the algebraic integers, hence an
algebraic integer, and all its conjugates have absolute
value $q^{1/2}$. Honda's theorem gives an abelian variety
$A/\mathbb F_q$ with eigenvalue $\beta$ on $H^1$.
The K\"unneth summand $H^1(A)^{\otimes i}$ of $H^i(A^i)$
then contains $\beta^i=\alpha$.
For $i=0$, Kronecker's theorem makes $\alpha$ a root of
unity of some order $m$. The permutation action on
$H^0(\operatorname{Spec}\mathbb F_{q^m}
\times_{\mathbb F_q}\overline{\mathbb F}_q,\mathbb Q_\ell)$
contains all $m$th roots of unity. This scheme is smooth
projective. If geometric connectedness were imposed,
the weight-zero assertion would instead allow only $1$.
\end{proof}
The weak assertion of 15/106 follows in this Weil-number
reading, even without a further extension of the finite
field. The claimed identification of the entire motivic
Galois group is separate: realization's full faithfulness
and algebraicity of invariant correspondences are not
proved merely by the existence of each eigenvalue.
Milne's Proposition 2.6 and Remark 2.7 retain Tate precisely
at that point.
More explicitly, his Conjecture 1.14 is the equality
between the dimension of numerical codimension-$r$
cycle classes and the order of the pole at $s=r$.
Proposition 1.15 derives algebraicity of all Frobenius
invariants and semisimplicity; Corollary 1.16 derives
full faithfulness of the $\ell$-adic realization.
These are the additional inputs used for the
injectivity part of Proposition 2.6 and for the
generation assertion of Remark 2.7. The eigenvalue
existence proof above supplies neither of them.
The integral condition also forces $i\geq0$:
for $i<0$ the nonzero integral norm of $\alpha$
would have absolute value less than one.

\paragraph{AQ-TATE-2026: the abelian sector is cleared in current literature.}
The hypothesis required for abelian varieties has now
been supplied in the primary literature. Van Hoften
and Sempliner (2026), Hexagon:2610.00259v2,
Theorem I, p.~1, prove the rational Hodge conjecture
for every complex CM abelian variety. Applying
Milne (1999), \emph{Lefschetz motives and the Tate
conjecture}, Theorem 7.1, journal p.~72, with the definition (0.1)
and finite-model descent in the Introduction, gives
the full Tate conjecture for every abelian variety
over every finite field. Milne's November 9, 2026
author note states this explicitly as Theorem 3.1,
p.~6. His 2002 paper's final author version,
Theorem 3.3(a),(b), p.~8, also gives numerical
equals $\ell$-adic homological equivalence and the
Hodge standard conjecture for abelian varieties
over algebraically closed fields; Remark 3.5,
p.~9, treats good Weil cohomology theories.
These are cited results, not independently reproved
in this module. The older long-version numbering
\emph{Theorem 4.6} used in the 2026 paper must not
be confused with Theorem 3.3 in the specified final
author PDF.

For comparison, Broe (2026), arXiv:2608.28651v1, Theorem 1.1,
printed p.~1, proves for every abelian fourfold
$A/\mathbb F_q$ and $\ell\ne\operatorname{char}\mathbb F_q$ that
\[
 CH^2(A)\otimes\mathbb Q_\ell\longrightarrow
 H^4_{\rm et}(A_{\overline{\mathbb F}_q},\mathbb Q_\ell(2))
 ^{\operatorname{Gal}(\overline{\mathbb F}_q/\mathbb F_q)}
\]
is surjective. With the divisor theorem and hard
Lefschetz this gives Tate in every codimension on $A$.
Corollary 1.2, p.~2, and Corollary 8.2, p.~20,
give agreement of numerical and $\ell$-adic homological
equivalence for abelian fourfolds over arbitrary fields
where $\ell$ is invertible. Corollary 8.1(i), p.~20,
also gives rational equals numerical equivalence with
rational coefficients over a finite field.
These are cited preprint results whose exact statements
were inspected here; their complete proofs are not
claimed independently checked. The fourfold result
alone does not address all powers of a fourfold;
the October--November all-abelian result does.
It therefore clears the realization issue on the
abelian--Artin category below. It does not prove
Tate for arbitrary smooth projective varieties,
or that every numerical motive belongs to that category.
Milne's Remark 7.6(b), journal p.~76, explicitly
records the extension to homological motives of
abelian type.

\paragraph{AQ-ABELIAN-RECOGNITION.}
Fix $k=\mathbb F_q$, $\ell\ne\operatorname{char}k$,
and let $\mathcal A_q$ be the rational numerical
category obtained from the motives of abelian
varieties over $k$, finite \'etale $k$-schemes, and
Tate motives by finite sums, tensor products,
duals, and direct summands. Equivalently, its
objects are direct summands of finite sums
\[
 h^i(A)(n)\otimes h^0(S),\qquad
 A/k\text{ abelian},\quad S/k\text{ finite \'etale},
 \quad n\in\mathbb Z.
\]
We use geometric Frobenius $F$ and
$\mathbb Q_\ell(1)$ with eigenvalue $q^{-1}$.
The summand above has weight $w=i-2n$ and
parity $i\bmod2=w\bmod2$.
Let $R_\ell(M)$ be its ordinary direct-sum
realization, with the weight decomposition remembered.
Then
\[
 \operatorname{Hom}_{\mathcal A_q}(M,N)\otimes_{\mathbb Q}\mathbb Q_\ell
 \xrightarrow{\ \sim\ }
 \operatorname{Hom}_{\mathbb Q_\ell[F]}
       (R_\ell(M),R_\ell(N)).                         \tag{AQ.15}
\]
Moreover, for $M,N\in\mathcal A_q$,
\[
 M\simeq N\quad\Longleftrightarrow\quad
 P_M(t)=P_N(t),\qquad
 P_M(t)=\det(1-tF\mid R_\ell(M)).                    \tag{AQ.16}
\]
The same equivalence holds with the weight-signed
Euler factor
\[
 Z_M(t)=\prod_w
 \det(1-tF\mid R_\ell(M)_w)^{(-1)^{w+1}}            \tag{AQ.17}
\]
in place of $P_M$. These statements are in the
rational numerical category, and concern actual
objects with their full multiplicities.

\begin{proof}
First we explain the scalar extension in (AQ.15).
For an abelian variety $X/k$, numerical and
$\ell$-adic homological equivalence agree by the
cited all-abelian theorem. Let $A^r_{\rm num}(X)_{\mathbb Q}$
denote the space of rational cycles modulo numerical
equivalence. The numerical intersection pairing
with $A^{\dim X-r}_{\rm num}(X)_{\mathbb Q}$ is nondegenerate.
For any finite linearly independent collection
$\alpha_1,\ldots,\alpha_s$ of numerical classes,
one can choose complementary classes
$\beta_1,\ldots,\beta_s$ making
$(\deg(\alpha_i\beta_j))_{i,j}$ invertible over
$\mathbb Q$. Their cohomology classes are then
linearly independent over $\mathbb Q_\ell$, by
pairing against the $\beta_j$. This also bounds
$s$ by the finite Betti number, proving finiteness
and injectivity after extension to $\mathbb Q_\ell$.
The full Tate theorem supplies surjectivity onto
the Frobenius-invariant subspace. Hence
\[
 A^r_{\rm num}(X)_{\mathbb Q}\otimes\mathbb Q_\ell
 \simeq H^{2r}(X_{\bar k},\mathbb Q_\ell(r))^{F=1}.
                                                               \tag{AQ.18}
\]

The cohomological projectors of an abelian variety
are algebraic here without a further conjecture.
For any integer $a\geq2$, multiplication by $a$
acts on $H^j(A)$ as $a^j$. The polynomials
\[
 e_i=\prod_{\substack{0\leq j\leq2\dim A\\j\ne i}}
       \frac{[a]^*-a^j}{a^i-a^j}
\]
give its projectors in numerical correspondences:
their required identities hold in cohomology,
hence numerically by the established equivalence
on $A\times A$. Tensor products are realized
using $A\times B$, and Poincar\'e duality expresses
duals as complementary cohomological degrees
with Tate twists. This proves the stated description
of $\mathcal A_q$.

Consider first $M=h^i(A)(n)$ and $N=h^j(B)(m)$,
and set $a=\dim A$. A correspondence between them
has codimension $a+m-n$ on $A\times B$. Its
relevant K\"unneth component, by Poincar\'e duality,
is
\[
 H^{2a-i}(A)\otimes H^j(B)(a+m-n)
 \simeq\operatorname{Hom}\bigl(H^i(A)(n),H^j(B)(m)\bigr).
\]
It lies in the required cohomological degree
precisely when $j-i=2(m-n)$, equivalently when
the weights agree. In that case (AQ.18) for
$A\times B$, followed by the algebraic projectors,
gives exactly (AQ.15). For unequal weights every
Frobenius-equivariant map is zero: the eigenvalues
have complex absolute values $q^{(i-2n)/2}$ and
$q^{(j-2m)/2}$, so their spectra are disjoint.
The numerical correspondence group is also zero,
since its cohomological component has the wrong
degree and homological equals numerical equivalence.

For finite \'etale factors $S,T$, the schemes
$A\times B\times S\times T$ are disjoint unions
of abelian varieties viewed over finite extensions
of $k$. A Frobenius-invariant class on a component
orbit is determined by a class invariant under the
appropriate power of Frobenius on one component.
The all-abelian Tate theorem over that finite
extension supplies the cycle. Intersection pairings
are multiplied by the nonzero extension degree,
which is invertible over $\mathbb Q$. Thus the
same injection and surjection apply to these
factors. Passing to finite sums and applying
source and target idempotents proves (AQ.15)
for every object of $\mathcal A_q$.

Frobenius acts semisimply on all these realizations.
For $H^1(A)$ this follows from Tate (1966),
Theorem 2(a), printed p.~140: $\mathbb Q[\pi_A]$
is the centre of the semisimple algebra
$\operatorname{End}_k(A)\otimes\mathbb Q$.
It is a product of number fields, so its Frobenius
element has square-free minimal polynomial;
dualizing or inverting preserves this property.
The identity $H^i(A)=\bigwedge^i H^1(A)$ shows
semisimplicity in every degree. Finite \'etale
Frobenius has finite order, and tensor products,
Tate twists and invariant direct summands of
semisimple single-operator representations remain
semisimple in characteristic zero.

Equality in (AQ.16) therefore gives an isomorphism
$R_\ell(M)\simeq R_\ell(N)$ as
$\mathbb Q_\ell[F]$-modules: each irreducible
polynomial in a semisimple operator's characteristic
polynomial records the multiplicity of its simple
$\mathbb Q_\ell[F]$-module. By (AQ.15) this is
an invertible element of
$\operatorname{Hom}_{\mathcal A_q}(M,N)\otimes\mathbb Q_\ell$.
Choose a rational basis $f_1,\ldots,f_s$ of this
Hom space. The polynomial
\[
 D(x_1,\ldots,x_s)=
 \det R_\ell\!\left(\sum_\nu x_\nu f_\nu\right)
\]
over $\mathbb Q_\ell$ is not zero. Since
$\mathbb Q$ is infinite, some rational tuple
has $D\ne0$; this elementary fact follows by
induction on the number of variables, using the
finite number of roots of a nonzero polynomial
in one variable. It gives a rational morphism
$f:M\to N$ with invertible realization.
By full faithfulness its realization's inverse
comes from a morphism in
$\operatorname{Hom}_{\mathcal A_q}(N,M)\otimes\mathbb Q_\ell$.
The equations $gf=1_M$ and $fg=1_N$ form a
finite linear system over $\mathbb Q$ in the
coefficients of $g$. A rational linear system
solvable after extension to $\mathbb Q_\ell$
is solvable over $\mathbb Q$, by row reduction.
Thus $f$ has a rational inverse. The converse
in (AQ.16) is immediate.

Finally a reciprocal root contributed by weight
$w$ has complex absolute value $q^{-w/2}$.
Weights of different parity cannot contribute
the same root. Hence the numerator and denominator
of (AQ.17) are coprime, each with constant term
one. Equality of those rational functions
recovers their numerator and denominator separately,
and their product is $P_M(t)$. This reduces
(AQ.17) to (AQ.16) and completes the proof.
\end{proof}
No essential-surjectivity claim onto all Frobenius
representations is needed here. In particular,
this recognition theorem does not identify an
arbitrary smooth projective motive with an object
of $\mathcal A_q$, does not retain integral lattices,
and does not classify mixed extensions. It is a
genuine strengthened answer to the abelian sector
of cote 15, supplied by the new cycle-algebraicity
theorem rather than by analytic zero location.

\paragraph{AQ-RATIONAL-FIBRE: an explicit neutralization obstruction.}
Even this now-established abelian category need not
admit a fibre functor with values in rational vector
spaces. This is a concrete restriction on any reading
of the archive's motivic Galois group as a neutral
group over $\mathbb Q$. It is not a refutation of a
group defined after scalar extension or of a motivic
gerbe formulation.
\begin{proof}
Consider $E:y^2=x^3-x$ over $\mathbb F_3$, then
base-change it to $k=\mathbb F_9$. It is smooth:
the derivative of $x^3-x$ is $-1$ in characteristic
three. For each of $x=0,1,2$, the right side is
zero, so there is one affine point; with infinity,
$\#E(\mathbb F_3)=4$. Thus its $3$-Frobenius
polynomial on $H^1$ is $T^2+3$. The square of
Frobenius is $-3$ on $H^1$, and injectivity of
the Tate-module action of endomorphisms implies
that the $9$-Frobenius endomorphism is $[-3]$.
Tate (1966), Theorem 2(d), p.~140, identifies
\[
 D=\operatorname{End}_{\mathbb F_9}(E)\otimes\mathbb Q
   =D_{3,\infty},
\]
the quaternion division algebra ramified at $3$
and infinity. Alternatively, the main theorem
on p.~134 gives $\dim_{\mathbb Q}D=4$, since
Frobenius is scalar on a two-dimensional Tate
module; every nonzero endomorphism of an elliptic
curve is an isogeny, so this algebra is a division
algebra.

Put $M=h^1(E)$ in $\mathcal A_9$. Pullback
endomorphisms give an injective algebra map
$D^{\rm op}\to\operatorname{End}(M)$.
By (AQ.15), the target has rational dimension
four, since its scalar extension is
$M_2(\mathbb Q_\ell)$. This map is therefore
an isomorphism, so $\operatorname{End}(M)$ is
a quaternion division algebra.

Use the usual parity-corrected symmetry on
numerical motives, for which the ordinary
$\ell$-adic realization is a symmetric tensor
functor. The categorical trace of $1_M$ is
then $2$: its image is the ordinary trace of
the identity on $H^1(E,\mathbb Q_\ell)$,
and the scalar map $\mathbb Q\to\mathbb Q_\ell$
is injective. A rational fibre functor would
preserve this trace, hence send $M$ to a
two-dimensional rational vector space. Its
faithfulness would embed the four-dimensional
division algebra $\operatorname{End}(M)$ into
$M_2(\mathbb Q)$. Equal dimensions force this
embedding to be an isomorphism. That is
impossible because $M_2(\mathbb Q)$ has
nonzero noninvertible elements. This proves
the asserted obstruction.

It also exhibits an object restriction invisible
in a list of eigenvalues alone. Here
$P_M(t)=(1+3t)^2$. There cannot be an object
$L\in\mathcal A_9$ with $P_L(t)=1+3t$:
otherwise (AQ.16) would give $L\oplus L\simeq M$,
but the former has a nontrivial idempotent
endomorphism and the latter's division algebra
has none. The eigenvalue $-3$ occurs, as
AQ-WEIL-EXIST predicts, but its multiplicity
cannot be discarded in a rational motivic
classification.
\end{proof}

\paragraph{AQ-WEIL-NOTATION: the recovered $Q'$ and the norm-one subgroup.}
The official cote 15 PDF has now been inspected beyond the
supplied excerpts.  Physical page 94, printed footer 93,
opens with the explicit definition
\emph{``$Q'$: cl\^oture tot[alement] r\'eelle de $\mathbb Q$''};
the bracket expands the handwritten abbreviation.
Thus $Q'=\mathbb Q^{\rm tr}$, the field of totally real
algebraic numbers, and the $i$ in $Q'(i)$ means
$\sqrt{-1}$, distinct from the integral weight index.
Physical page 98/footer 97 discusses a totally imaginary
quadratic extension of a totally real field and writes
$\lambda\overline\lambda=q^w$, followed by the
normalization $\lambda=q^{w/2}u$ with
$u\overline u=1$.
Physical page 100/footer 99 contains the boxed correction
requiring every conjugate to have absolute value $q^{w/2}$,
as already preserved in the supplied transcription.
These recovered norm and purity conditions identify the
weight-zero interpretation used in AQ-WEIL-LIMIT.

More precisely, let
\[
 \mathbb Q^{\rm CM}=\bigcup_{L\ {\mathrm{CM}}}L
 \quad\hbox{inside the fixed }\overline{\mathbb Q}\subset\mathbb C.
\]
Then
\[
 \mathbb Q^{\rm CM}=\mathbb Q^{\rm tr}(\sqrt{-1}).
 \tag{AQ.WN.1}
\]
For a prime power $q=p^n$, the group of algebraic
$p$-units with all conjugate absolute values $q^{w/2}$,
for some integer $w$, is exactly
\[
 \left\{\lambda\in\mathbb Q^{\rm tr}(\sqrt{-1})^*:
   \lambda\text{ is a $p$-unit},\quad
   \lambda\overline\lambda=q^w\text{ for some }w\in\mathbb Z
 \right\}.                                                   \tag{AQ.WN.2}
\]
Its weight-zero subgroup is $W_0(p)$, independent of $n$.
Thus the historical group $M(q)$ and its weight-zero part
agree with $W(q)$ and $W_0(p)$ when the explicit
norm/purity conditions on pages 98 and 100 are retained.

\begin{proof}
For a CM field $L$, let $L^+$ be its totally real
subfield.  Choose $b\in L$ with
$\delta=b-\overline b\ne0$.  Then
$\delta^2\in L^+$, and $a=-\delta^2$ is totally positive:
every embedding sends $\delta$ to a nonzero purely
imaginary number.  The algebraic number $\sqrt a$
is totally real, because all its conjugates are real
square roots of positive conjugates of $a$.
Hence
$L=L^+(\delta)\subset\mathbb Q^{\rm tr}(\sqrt{-1})$.
Conversely any element of the latter field belongs to
$F(\sqrt{-1})$ for a finite totally real field $F$,
and $F(\sqrt{-1})$ is a CM field.  This proves (AQ.WN.1),
including the assertion that the displayed union is a field.

Let $\lambda$ have the absolute values in (AQ.WN.2),
and let $L$ be the Galois closure of $\mathbb Q(\lambda)$.
Write $c$ for complex conjugation on $L$.
For every $\sigma\in\operatorname{Gal}(L/\mathbb Q)$,
\[
 c(\sigma\lambda)=q^w/(\sigma\lambda).
\]
For every $\tau$ in the same Galois group,
$c\tau$ and $\tau c$ consequently agree on all the
conjugates $\sigma\lambda$, which generate $L$.
Thus $c$ is central.
If $c=1$, the Galois field $L$ is totally real.
Otherwise $L^+=L^{\langle c\rangle}$ is totally real,
$[L:L^+]=2$, and $L$ has no real embedding, so $L$ is CM.
In either case (AQ.WN.1) places $\lambda$ in
$\mathbb Q^{\rm tr}(\sqrt{-1})$; the norm identity follows
from its specified absolute value at the fixed embedding.

Conversely this field is Galois over $\mathbb Q$, and
complex conjugation fixes $\mathbb Q^{\rm tr}$ and
negates $\sqrt{-1}$.  It therefore commutes with every
embedding.  Applying any embedding to
$\lambda\overline\lambda=q^w$ gives
$|\sigma\lambda|^2=q^w$, as required.
The $p$-unit condition is the same on both sides,
so (AQ.WN.2) follows.  Setting $w=0$ gives the claimed
weight-zero subgroup, independently of $q=p^n$.
\end{proof}

Finite-prime valuations alone must not be substituted
for the norm/purity condition in this identification.
For example $\varepsilon=2+\sqrt3$ and its inverse
$2-\sqrt3$ are algebraic integers, so $\varepsilon$
is a unit at every finite place.  In particular every
finite-prime sum $v_{\mathfrak p}(\varepsilon)+
v_{\overline{\mathfrak p}}(\varepsilon)$ is zero,
for every chosen $p$.  But
$\varepsilon\overline\varepsilon=\varepsilon^2\ne1$
and its two real conjugates have unequal absolute values.
No positive power is a weight-zero Weil number.
This checks the limitation of a \emph{valuation-only}
paraphrase of the margin on 15/102; it does not assert
that the manuscript's terminology \emph{unit\'e absolue}
itself discards archimedean conditions.
The corrected group interpretation uses the explicit
norm condition actually present on the preceding sheet.

\paragraph{AQ-WEIL-LIMIT: the erased question at 15/104.}
The text \emph{Trouve-t-on $P^0/U$ tout entier?}
and its subsequent root-extraction test are inside a
deletion in the supplied transcription. In the standard
Weil-number interpretation its algebraic assertion is true.
Let $W_0(p)$ be the multiplicative group of algebraic
$p$-units whose complex conjugates all have absolute
value one. Then the directed system with transition
maps $u\mapsto u^m$ has limit
\[
 \varinjlim W_0(p)=W_0(p)\otimes_{\mathbb Z}\mathbb Q
                 \simeq W_0(p)/\mu_\infty .
\]
\begin{proof}
For any abelian group $A$, the system indexed by
positive integers, with multiplication by $m$ from
index $n$ to index $mn$, has colimit
$A\otimes\mathbb Q$: the compatible map sends
the class of $a$ at index $n$ to $a/n$.
Every root of an element of $W_0(p)$ is again
a $p$-unit with all conjugate absolute values one.
Thus $W_0(p)$ is divisible, and its torsion subgroup
is exactly $\mu_\infty$. Its quotient by torsion
is uniquely divisible, giving the last isomorphism.
\end{proof}
The recovered original definition $Q'=\mathbb Q^{\rm tr}$
and the preceding explicit norm-one condition identify
this weight-zero interpretation by AQ-WEIL-NOTATION.
The literal deleted calculation therefore has the stated
affirmative answer with its norm/purity conditions retained.
Its identification with the motivic group of arbitrary
motives still requires the separate Tate, generation,
fibre and base-field transition conventions described above;
the abelian realization sector has been cleared by
AQ-TATE-2026 and CC-AG01--03.

\paragraph{Coefficient-extended group: the completed abelian-sector answer.}
The companion proof CC-AG01--03, in
\texttt{research/consequences\_criteria\_abelian\_group.tex},
now strengthens the rational recognition result to a full
category and group description. Fix
$E=\overline{\mathbb Q}_\ell$ and an embedding
$\overline{\mathbb Q}\hookrightarrow E$. Let $W(q)$ be the
group of algebraic $p$-units whose conjugate absolute values
are $q^{w/2}$ for one integer $w$. With the parity-corrected
symmetry and idempotent completion after extending Hom spaces,
the result is
\[
 \operatorname{Kar}(\mathcal A_q\otimes_{\mathbb Q}E)
  \simeq \operatorname{Gr}^{\mathrm{fin}}_{W(q)}(E),\qquad
 \operatorname{Aut}^{\otimes}(\omega_E)=D_E(W(q)).
\]
The target means finite-dimensional vector spaces graded by
$W(q)$, with multiplication of labels as tensor product;
$D_E(M)=\operatorname{Spec}E[M]$ for an abelian group $M$.
Choosing $\sqrt q$ gives
$D_E(W(q))\simeq D_E(W_0(p))\times\mathbb G_{m,E}$,
with Tate character $(1,-2)$. The full proof uses the
Hom-space theorem proved above, Honda occurrence, and
lifting of the individual eigenspace projectors over $E$.
It also proves the distinguished Frobenius has Zariski-dense
powers. The finite-field group is not a protorus, since its
character group contains $\mu_\infty$. This answers the
abelian-sector group question with explicit coefficient and
fibre conventions, while preserving the rational quaternion
obstruction and distinguishing passage to
$\overline{\mathbb F}_p$ from extension of coefficients.
It does not assert that all numerical motives are of
abelian type or identify the fixed-finite-field category
with the separate category of Frobenius germs over
$\overline{\mathbb F}_p$.  The previously missing field
notation is now identified by AQ-WEIL-NOTATION.

\subsection{AQ-TWIST: the weight-index correction at 18/009}
The sheet's conventions give
\[
 W_j(V(n))=W_{j+2n}(V),\quad F^j(V(n))=F^{j+n}(V),\quad
 \operatorname{gr}^W_j(V(n))
 =(\operatorname{gr}^W_{j+2n}V)(n).                             \tag{AQ.6}
\]
The last equality follows by taking the quotient of the
first two consecutive weight steps, with the shifted Hodge
filtration. The transcribed unshifted graded formula fails
for $V=\mathbb Q(0)$, $n=-1$, $j=2$: the left side is
$\mathbb Q(-1)$ and the unshifted right side is zero.

\subsection{AQ-REAL: the even-Tate assertion at 18/038 is false}
The sheet claims that $f_{\infty,M}=1$ forces a sum of
tensor powers of $\mathbb Q(2)$.
Let $A$ be the sign Artin motive over $\mathbb R$ and
$L=A\otimes\mathbb Q(-1)$. Its real involution is
$(-1)(-1)=1$, but its weight $2$ is not the weight $-4r$
of any $\mathbb Q(2)^{\otimes r}$. This refutes the claim.
In the semisimple real Artin--Tate category the correct
objects are exactly finite sums of $L^{\otimes r}$,
$r\in\mathbb Z$.
\begin{proof}
The simple objects are $A^\epsilon\otimes\mathbb Q(n)$,
$\epsilon\in\{0,1\}$, with real involution
$(-1)^{\epsilon+n}$. It is trivial exactly when
$\epsilon\equiv n\pmod2$, in which case this object is
$L^{\otimes(-n)}$. Direct-sum decomposition proves both
directions.
\end{proof}
The classification is asserted here only for Artin--Tate
motives. Extending it from Hodge types to all pure motives
would require an additional algebraicity input.

\subsection{AQ-PERIOD-STRONG: a counterexample to 18/011}
The strong transcendence conjecture on 18/011 predicts
$\operatorname{trdeg}_k k(\alpha,\beta)=g+u$.
In the usual reading $T_{\rm Hdg}=\operatorname{gr}_F T_{\rm dR}$
and $Q$ parametrizes normalized splittings, this fails.
Take $E:y^2=x^3-x+1$ over $\mathbb Q$.
Its $j=-6912/23$ is not integral, so it is non-CM;
for $M=h^1(E)$ one has $G=\mathrm{GL}_2$, $g=4$, $u=1$.
Here we use the classical integrality of CM
$j$-invariants, recalled in Habegger (2014),
\S1, p.~1, and the elliptic motivic-group
identification in Kreutz--Shen--Vial,
Proposition 5.14.
Here $\alpha$ is the comparison from Betti to de Rham
cohomology, $T_{\rm Hdg}=\operatorname{gr}_F T_{\rm dR}$,
$\beta$ is the Hodge splitting normalized to induce the
identity on that associated graded, and
$\gamma=(\beta\alpha)^{-1}$, as in the source triangle.
The parabolic preserving its Hodge line has
one-dimensional unipotent radical, explaining $u=1$.
\begin{proof}
Choose rational Betti eigenvectors for the real involution
and a rational de Rham basis with Hodge line $\mathbb Ce_1$.
Write
\[
 P=\begin{pmatrix}a&b\\c&d\end{pmatrix},\quad
 \delta=ad-bc\ne0,\quad C=\operatorname{diag}(1,-1).
\]
The opposite Hodge line is $PCP^{-1}\mathbb Ce_1$,
generated by $(ad+bc,2cd)^t$. Both $c,d$ are nonzero:
otherwise a real-involution eigenline would be the
Hodge line, forcing it to equal its opposite.
Hence the normalized splitting is
\[
 \beta=\begin{pmatrix}1&-(ad+bc)/(2cd)\\0&1\end{pmatrix}.
\]
It is rational in $P$, so $k(P,\beta)=k(P)$ has
transcendence degree at most four, not $g+u=5$.
Moreover
\[
 \gamma=(\beta P)^{-1}
 =\begin{pmatrix}d/\delta&1/(2c)\\-c/\delta&1/(2d)\end{pmatrix},
 \qquad \gamma_{11}\gamma_{22}+\gamma_{12}\gamma_{21}=0.
\]
Thus its transcendence degree is at most three,
also contradicting the predicted value $g=4$.
\end{proof}
This refutation does not settle the ordinary non-CM
elliptic period conjecture or algebraicity of its
Hodge-splitting scalar.

\subsection{AQ-PERIOD-GENERATION: 18/013(b) is a known theorem}
Question (a) asks for period-field transcendence degree
$\dim G$, namely four without CM or two with CM.
Question (b) asks whether the group generated by a
translate of the period locus is all of $G$.
The answer to (b) for one elliptic curve is yes:
Kreutz, Shen, and Vial (2026), \emph{De Rham--Betti classes
with coefficients}, arXiv:2206.08618v3, Proposition 3.1
and Theorem 5.13(ii) identify the generated period torsor
with the motivic torsor; Proposition 5.14 identifies
the elliptic motivic group. The coefficient restriction
in Theorem 5.13(ii) is satisfied for a single elliptic
curve. This known result is cited, not reproved.
It is weaker than density of the period locus.

Question (a) is known for CM curves by Chudnovsky (1980),
\S2A, Theorem 1 and Corollary 2, printed p.~343
in the 1978 Helsinki ICM proceedings. In particular,
a nonzero CM period and $\pi$ are algebraically
independent. The non-CM case remains the full
period-conjecture gate. Question (c) asks whether the
Hodge-splitting scalar can be algebraic without CM.
In the normalization $y^2=4x^3-g_2x-g_3$,
put $\omega=dx/y$, $\xi=x\,dx/y$, take full periods
$\omega_1,\omega_2$ with $\operatorname{Im}
(\omega_2/\omega_1)>0$, and write the periods of
$\xi$ as $-\eta_1,-\eta_2$. The opposite Hodge line
is generated by $\xi+h\omega$. Its periods must be
proportional to $\overline{\omega_1},
\overline{\omega_2}$, which gives
\[
 h=\frac{\eta_1\overline{\omega_2}-\eta_2\overline{\omega_1}}
          {\omega_1\overline{\omega_2}-\omega_2\overline{\omega_1}}
   =\frac{\pi^2}{3\omega_1^2}E_2^*(\tau).
\]
Here $\tau=\omega_2/\omega_1$ and
$E_2^*(\tau)=E_2(\tau)-3/(\pi\operatorname{Im}\tau)$.
The second equality uses
$\eta_1\omega_2-\eta_2\omega_1=2\pi i$ and
$\eta_1\omega_1=\pi^2E_2(\tau)/3$.
The CM conclusions in (a) and (c) can be separated
from the remaining non-CM problem.
\begin{proof}[The CM period field and algebraicity of the splitting]
Take a nonreal CM endomorphism $u$, defined over
$\overline{\mathbb Q}$. In the algebraic de Rham
basis its action has the form
$u^*\omega=u\omega$ and
$u^*\xi=a\omega+\overline u\xi$, with
$a\in\overline{\mathbb Q}$. The opposite Hodge line
is the $\overline u$-eigenline, so
\[
 a+hu=\overline u h,\qquad
 h=\frac{a}{\overline u-u}\in\overline{\mathbb Q}.
\]
The period ratio $\tau$ is imaginary quadratic.
The preceding formula for $h$ and the Legendre
relation give
\[
 \omega_2=\tau\omega_1,\qquad
 \eta_1=h\omega_1+\frac{\pi}{\omega_1\operatorname{Im}\tau},
 \qquad \eta_2=\tau\eta_1-\frac{2\pi i}{\omega_1}.
\]
Thus the full period and quasi-period field is
$\overline{\mathbb Q}(\omega_1,\pi)$; the reverse
inclusion follows from the Legendre relation.
Chudnovsky's cited Corollary 2 makes its
transcendence degree exactly two.
\end{proof}
The remaining non-CM assertion is the nonholomorphic weight-two
Eisenstein-transcendence question attributed to Katz
in Waldschmidt (2004), p.~277, immediately after
Conjecture 3.22. That passage states the same remaining
condition for a non-CM lattice with algebraic invariants.
Its non-CM transcendence is not established by the
supplied proofs. AQ-FINITE cannot supply that missing
algebraic-independence assertion.

\paragraph{AQ-PERIOD-2026: current-source and normalization check.}
The current primary literature does not identify the
three different weight-two functions in these questions.
Kreutz--Shen--Vial, arXiv:2206.08618v3, Introduction
and \S4.2, retain the full period conjecture separately
from their generated-torsor theorem. Mohajer (2026),
arXiv:2611.25399v1, \S1.1, equation (1.3), pp.~2--3,
still formulates the non-CM complex four-period
transcendence-degree assertion as conjectural.
His Theorems 4.2 and 4.4 concern comparison periods
in $B_{\rm dR}$ and their normalized images in
$\mathbb C_p$, respectively. His Proposition 2.1
concerns zeros of the holomorphic $E_2$, whereas
the scalar above involves the almost holomorphic
$E_2^*=E_2-3/(\pi\operatorname{Im}\tau)$.

Sebbar and Besrour (2026), arXiv:2611.14502v1,
\S8.4, Remark 8.11, Proposition 8.12 and
Question 8.13, pp.~32--33, consider the ordinary
$p$-adic Katz form $E^{\rm Katz}_2$. They leave
the existence of a non-CM zero with algebraic
$j$-invariant over $\mathbb Q$ as a question;
their nonvanishing proposition assumes the crystalline
period conjecture of Ancona--Fr\u a\c til\u a.
Algebraicity over a finite extension of $\mathbb Q_p$
does not answer algebraicity over $\mathbb Q$.
Thus these November 2026 results do not supply the
missing complex algebraic-independence assertion in
18/013(a), or the non-CM algebraicity classification of
the complex splitting scalar in 18/013(c). This is a
source-scope conclusion, not a proof of logical
nonimplication from the adopted analytic theorems.

\paragraph{AQ-CM-PERIOD-2026: the new rank-one comparison results.}
Kahn and Otsubo (2026), arXiv:2611.13723v1,
Corollary 3.21, p.~24, show that equal Hecke characters
give equal complex period classes for rank-one abelian-type
$K$-coefficient motives over a number field $k\supset K^{\rm gal}$.
Here the class lies in
$(\mathbb C\otimes_{\mathbb Q}K)^*/(k\otimes_{\mathbb Q}K)^*$
(\S2.4, pp.~6--7).
Remark 3.22, p.~25, relates the converse, up to a finite-order
character, to Rohrlich's conjecture.
For rank-one algebraic motives, equality of crystalline period
classes at almost all places implies equality of Hecke characters;
the converse holds under the preceding abelian-type and base-field
hypotheses (Corollary 3.35, pp.~30--31).
Corollary 3.25, p.~26, proves Mumford--Tate when the rank-one
$K$-motive becomes algebraic after finite extension.
For an abelian CM field $K$ of conductor $m$ and
$k\supset\mathbb Q(\mu_m)$,
Theorems 4.36 and 4.39, pp.~46--47,
refine Gross--Deligne, and Theorem 4.44, p.~49,
gives the corresponding $p$-adic beta-period formula.
These are positive comparison results with specified hypotheses.

For a non-CM elliptic curve, $h^1(E)$ has rational rank two
and endomorphism algebra $\mathbb Q$.  A rank-one structure
over a coefficient field would require a quadratic field to
embed in that endomorphism algebra, which is impossible.
The determinant is a Tate object and records only the
determinant period.  Thus these rank-one statements do not
give the four-period algebraic independence or the
complex $E_2^*$ assertion in 18/013.

\paragraph{AQ-KO-RANK: a correction to an ancillary dimension claim.}
With Definition 2.5, p.~8, of the same paper, the dimension
sentence in Lemma 4.6, p.~34, requires one additional rank:
for $K=\mathbb Q(\mu_m)$ with $m>2$,
\[
 \operatorname{rank}_{\mathbb Z}\Xi(K)
       =\frac{\varphi(m)}2+1,
 \tag{AQ.KO.1}
\]
where the paper prints $\varphi(m)/2$.
\begin{proof}
Put $r=\varphi(m)/2$ and choose one embedding
$\sigma_j$ from each of the $r$ complex-conjugate pairs.
Since $K$ itself is a CM field, the additional CM-subfield
condition in that definition is automatic.
A function $\phi:\Sigma_K\to\mathbb Z$ belongs to
$\Xi(K)$ precisely when all the sums
$\phi(\sigma_j)+\phi(\overline\sigma_j)$ equal
one common integer $w$.
Consequently
\[
 \Xi(K)\longrightarrow\mathbb Z^{r+1},\qquad
 \phi\longmapsto(w,\phi(\sigma_1),\ldots,\phi(\sigma_r))
\]
is an isomorphism of abelian groups: the inverse sends
$(w,a_1,\ldots,a_r)$ to the function with values
$a_j$ and $w-a_j$ on the $j$th pair.
This proves (AQ.KO.1).  The map
$\phi\mapsto2\phi-w$ used in the cited rank argument
has the full group of constant functions as its kernel,
rather than being injective.  Already $m=3$ gives
$\Xi(\mathbb Q(\mu_3))=\mathbb Z^2$.
\end{proof}
This calculation concerns that rank sentence; the
period-comparison results above are cited with their
own hypotheses, and their complete external proofs
are not certified by this source-scope check.

\subsection{AQ-COCYCLE: the image-density statement at 18/059}
Let $P$ be an affine torsor under an algebraic group $G$
over $F\subset\mathbb C$, and $p\in P(\mathbb C)$.
The closure of the image of its automorphism cocycle
is the translated $F$-locus of $p$. Thus the image is
Zariski dense in $G_{\mathbb C}$ exactly when $p$ is
generic in $P$ over $F$.
\begin{proof}
Embed $P$ in affine space over $F$, and let $I$ be the
ideal over $\mathbb C$ of the
$\operatorname{Aut}(\mathbb C/F)$-orbit of $p$.
Each finite-dimensional degree-bounded part of $I$
is stable under the coefficientwise semilinear action.
Its unique reduced-row-echelon basis, relative to an
ordered monomial basis, must therefore have entries
fixed by every automorphism. The fixed field of
$\operatorname{Aut}(\mathbb C/F)$ is $F$: any nontrivial
embedding of $F(a)$ over $F$, for $a\notin F$, extends
first across a transcendence basis and then an algebraic
closure to an automorphism of $\mathbb C$.
Consequently every degree-bounded part has a basis
over $F$, and $I=(I\cap F[\mathbf x])\mathbb C[\mathbf x]$.
An $F$-polynomial vanishes on the orbit exactly when
it vanishes at $p$. Thus $I$ is precisely the ideal of
$Z_F(p)_{\mathbb C}$, proving the density assertion.
The torsor difference map $D_p:P_{\mathbb C}\to G_{\mathbb C}$
is an algebraic isomorphism and the cocycle is
$D_p(\sigma p)$, up to inversion according to convention.
Its closure is consequently $D_p(Z_F(p)_{\mathbb C})$.
\end{proof}
The generated-subgroup version is false for general
torsors. In $G=\mathbb G_m^2$ take $F=\mathbb Q$ and
$p=(t,t+1)$ with $t$ transcendental. Its locus is a curve.
The translated cocycle curve is
$u\mapsto(u/t,(u+1)/(t+1))$.
A character $x^a y^b$ identically one there would have
orders zero at $u=0$ and $u=-1$, giving $a=b=0$.
Thus this proper curve generates the whole torus.
This refutes an abstract torsor inference, without
claiming an arithmetic-motivic counterexample to the
period conjecture.

\subsection{AQ-CM: the hexagonal elliptic curve and a genuine GRH consequence}
The hexagonal real-elliptic classification recorded in SA22,
\S2(A.7), concerns the $j=0$ period lattice and its real
forms. It is not the Grothendieck--Teichm\"uller hexagon.
For example, the latter is the identity
\[
 f(z,x)z^m f(y,z)y^m f(x,y)x^m=1,\qquad
 z=(xy)^{-1},\quad m=(\lambda-1)/2,
\]
in a profinite free group, as in Boavida de Brito, Horel,
and Robertson (2019), Definition 7.2(II). The two uses of
\emph{hexagonal} name different mathematical data.

\paragraph{Hodge obstruction.}
For $E:y^2=x^3+1$, the motive $h^1(E)$ is not Artin--Tate,
even after adjoining its CM field $\mathbb Q(\sqrt{-3})$.
\begin{proof}
An Artin motive has Hodge type $(0,0)$, and $\mathbb Q(n)$
has type $(-n,-n)$ in our convention.
Tensor products, sums, subobjects, and quotients of such
Hodge structures have only diagonal types. But
$H^1(E)$ has nonzero types $(1,0)$ and $(0,1)$.
Base extension to an algebraic number field does not
change these complex Hodge types. This proves the claim.
\end{proof}
Thus the fact that the CM field is abelian does not turn
its weight-one Hecke character into a finite-order
Dirichlet character with an integral Tate twist.
No conclusion about the zeros of $L(h^1(E),s)$ is obtained
merely by substituting that field in the supplied
Dirichlet theorem.

\paragraph{Effective distribution of trace-zero primes.}
Let $a_p=p+1-\#E(\mathbb F_p)$, for $p>3$, and put
$N_0(x)=\#\{3<p\leq x:a_p=0\}$. Then
\[
 a_p=0\ \Longleftrightarrow\ p\equiv2\pmod3,\qquad
 N_0(x)=\tfrac12\operatorname{Li}(x)+O(\sqrt x\log x)            \tag{AQ.7}
\]
under the adopted Dirichlet GRH. The implied constant is
effective.
\begin{proof}
With the quadratic character extended by $\chi_p(0)=0$,
\[
 a_p=-\sum_{x\in\mathbb F_p}\chi_p(x^3+1).
\]
If $p\equiv2\bmod3$, cubing permutes $\mathbb F_p$,
so the sum is $\sum_u\chi_p(u+1)=0$.
If $p\equiv1\bmod3$, set $m=(p-1)/2$.
Euler's criterion and the binomial theorem show
\[
 a_p\equiv-\sum_{j=0}^m {m\choose j}
                 \sum_{x\in\mathbb F_p}x^{3j}
       ={m\choose (p-1)/3}\pmod p.
\]
The $j=0$ sum is $p=0$ modulo $p$.
Among the positive exponents $3j\leq3(p-1)/2$, the
only multiple of $p-1$ is $p-1$, and its power sum is
$-1$ modulo $p$. The displayed binomial coefficient
is nonzero modulo $p$ because all its factorial arguments
are less than $p$. This proves the equivalence.

Now $\zeta_{\mathbb Q(\sqrt{-3})}(s)=
\zeta(s)L(s,\chi_{-3})$, so the adopted theorem gives
GRH for this quadratic Dedekind zeta function.
Lagarias--Odlyzko (1977), Theorem 1.1, equation (1.3),
applied to its nontrivial Frobenius class, gives
\[
 \left|\pi(x;3,2)-\tfrac12\operatorname{Li}(x)\right|
 \leq c_1\left(\tfrac12\sqrt x\log(3x^2)+\log3\right),
 \qquad x>2.
\]
Here $c_1$ is their effective absolute constant.
For $x\geq3$, $N_0(x)=\pi(x;3,2)-1$ because $p=2$
is a bad prime of $E$. This proves (AQ.7).
\end{proof}
This gives an effective arithmetic statement about the
hexagonal curve using the quadratic splitting character.
It does not replace that curve's weight-one Hecke
$L$-function by a Dirichlet $L$-function.

\subsection{AQ-67-RANGE: the global isomorphism in 67/078 is false}
The literal question is \emph{Ces fl\`eches sont elles
des iso --- et si oui, comment expliciter un iso en
sens inverse ?} In the stable nodal $(0,4)$ case,
normalize both three-pointed rational components with
the node at $0$ and the other two marked points at
$1,\infty$. In the tangent-tensor normalization of
the capping construction, a tensor $q$ in the image
has absolute value equal to the product of the two
node-cap conformal radii. Consequently
\[
 |q|\leq16,\qquad q=17\ \hbox{has no preimage}.                  \tag{AQ.8}
\]
If the target is written with dual tangent tensors,
the equivalent statement uses the inverse tensor.
This normalization is fixed by the original sheets:
67/075 calls the propagation datum a nonzero basis of
$T_{x'}\otimes T_{x''}$; 67/078 writes its inverse
as a dual pairing and asks for mutually polar disks
in the two tangent lines. In coordinates, polarity
is $|q^{-1}|r'r''=1$, hence $r'r''=|q|$.
The target in 67/076--077 imposes no bound on this
nonzero tensor. We use the map without additional
marked-point disks in 67/077.

\begin{proof}
For completeness, the Koebe bound needed here follows
from the area argument. If
$F(z)=z+b_0+\sum_{n\geq1}b_nz^{-n}$ is univalent for
$|z|>1$, Green's formula gives the area inside
$F(|z|=r)$ as
$\pi(r^2-\sum_{n\geq1}n|b_n|^2r^{-2n})\geq0$.
Letting $r\downarrow1$ proves
$\sum_{n\geq1}n|b_n|^2\leq1$.
For a univalent $f(z)=z+a_2z^2+\cdots$ in the unit disk,
$g(z)=\sqrt{f(z^2)}=z+(a_2/2)z^3+\cdots$ is an odd
univalent function there. The branch exists because
$f(z^2)/z^2$ never vanishes in the disk. Equality of
two $g$-values forces equality of the squared
arguments under $f$, and oddness then proves
injectivity. Its reciprocal
$1/g(1/z)$ is univalent outside the disk, so the area
inequality gives $|a_2|\leq2$.
If $w\ne0$ is omitted by $f$, the normalized univalent
function $f/(1-f/w)$ has second coefficient $a_2+1/w$.
Thus $|1/w|\leq|a_2|+|a_2+1/w|\leq4$.
This is the Koebe quarter theorem (Duren, 1983,
Theorem 2.3, p.~31).

A node cap is the image of a univalent map $h$ from
the unit disk with $h(0)=0$ and conformal radius
$r=|h'(0)|$. It omits $1$ and $\infty$.
Apply the preceding conclusion to $h/h'(0)$ and
the omitted value $1/h'(0)$ to get $r\leq4$.
Both caps therefore have radii at most four.
The tensor is the product of the two derivative
vectors, up to the unit-modulus boundary identification,
so its norm is at most sixteen. The nonzero tensor
$17$ belongs to the stated unrestricted target,
establishing failure of surjectivity.
\end{proof}
The intake SA22 reports the sharper exact image
$0<|q|<16$, tangent invertibility, and a $\pi_1$-isomorphism.
Its full prior proofs are not in the supplied packet.
The next result independently supplies complete proofs
for the labelled geodesic-cut $(0,4)$ sector, while
AQ-67-NONHOLO separates the real and complex tangent claims.
The elementary obstruction (AQ.8) already answers
67/078 negatively and is independent of RH.

\subsection{AQ-67-EXACT: the canonical radius and the exact labelled $(0,4)$ image}
\paragraph{Scope and normalization.}
Consider the geodesic-cut interpretation of the map without marked-point
disks in cote 67/077. Fix a partition of four labelled punctures into two
pairs and remember the separating geodesic. Its length is $\ell>0$.
Cutting produces two hyperbolic pairs of pants of boundary lengths
$(\ell,0,0)$, where zero denotes a cusp. Cap the geodesic boundary by the
unit disk using its normalized hyperbolic arclength, as prescribed in
67/074. On each resulting three-pointed sphere use the unique coordinate
with cap centre $0$ and the two labelled cusps $1,\infty$.
Let $r(\ell)>0$ be the conformal radius of the cap at zero in that
coordinate. A cap coordinate $w$ has $|w|=1$ at the seam and
$z=r(\ell)w+O(w^2)$ after a rotation of $w$.
The tensor coordinate $q$ of 67/075 then has modulus $r(\ell)^2$.

Put
\[
 x=\frac{\ell}{4\pi},\qquad
 \phi(x)=\operatorname{Im}\bigl(\log\Gamma(1+ix)
                       -\log\Gamma(\tfrac12+ix)\bigr),
\]
where the analytic logarithm of $\Gamma$ on the right half-plane is
real on the positive real axis. Then
\[
 \boxed{\quad r(\ell)=4\exp\!\left(\frac{\phi(x)-\pi/4}{x}\right).\quad}
                                                               \tag{AQ.67.1}
\]
The function $r:(0,\infty)\longrightarrow(0,4)$ is a real-analytic
diffeomorphism. With the orientation of the normalized twist angle
$\theta\in\mathbb R/2\pi\mathbb Z$ chosen so that its increase multiplies
the tangent tensor by $e^{i\theta}$, the capping map is
\[
 (\ell,\theta)\longmapsto r(\ell)^2e^{i\theta}.
                                                               \tag{AQ.67.2}
\]
It is a real-analytic diffeomorphism onto $0<|q|<16$.
Its real tangent map is invertible everywhere, and its composite with
the inclusion $\{0<|q|<16\}\hookrightarrow\mathbb C^\times$ induces an
isomorphism of fundamental groups. These are statements about this
labelled geodesic-cut sector. They do not assert complex holomorphicity
or algebraic \emph{\'etaleness} of the map in 67/077--082, nor do they
settle the full multi-node question.

\paragraph{Geometric inputs and source locators.}
The hyperbolic pair of pants with specified nonnegative boundary
lengths is unique, with zero lengths interpreted as cusps. Its
Fenchel--Nielsen length and twist parametrization is recalled in
Mondello (2011), \S2.2, p.~640. The canonical attachment here is the
infinite flat-cylinder grafting of Mondello, \S5.4,
Theorem 5.4, p.~660; a punctured flat cylinder is conformally the
punctured disk with the arclength parametrization used above.
For an explicit independent source for the projective-connection
calculation below, see F\i rat (2021), equations (3.11), (3.13),
(3.14a), (3.23), and (3.30)--(3.34), arXiv:2102.03936v2,
pp.~14--17 and 24--25. We specialize the two other ends to cusps
and determine the branch directly below; we do not infer the branch
from that paper's integer-valued numerical prescription (3.27).

\begin{proof}
Write $\lambda=\ell/(2\pi)$, so $x=\lambda/2$.
The developing map $D$ of the hyperbolic metric is a local isometry
to a half-plane. Near a geodesic of length $2\pi\lambda$, diagonalize
its hyperbolic monodromy, with fixed points $0,\infty$.
In the arclength cap coordinate its analytic continuation is, up
to a nonzero constant, $w^{i\lambda}$. Indeed, increasing
$\arg w$ by $2\pi$ multiplies this function by $e^{-2\pi\lambda}$,
and along $|w|=1$ its hyperbolic arclength is
$\lambda\,d\arg w$. The same expression supplies the analytic
continuation across the seam. Its Schwarzian derivative therefore
has a double pole at the cap centre with leading coefficient
$(1+\lambda^2)/2$. At a cusp its leading coefficient is $1/2$,
as follows from the local developing coordinate $\log(z-z_0)$.
There are no other poles. On the sphere with centres and cusps
$0,1,\infty$, the resulting Schwarzian is consequently
\[
 T(z)=\frac{1+\lambda^2}{2z^2}
       +\frac1{2(z-1)^2}
       -\frac{1+\lambda^2}{2z(z-1)}.                           \tag{AQ.67.3}
\]
To see uniqueness, any two such projective connections differ by
a quadratic differential with at most simple poles at three points;
this space is $H^0(\mathbb P^1,\mathcal O(-1))=0$.

Set $a=(1+i\lambda)/2$. Direct substitution in
$y''+\tfrac12T(z)y=0$ gives the two solutions
\[
 y_\pm(z)=z^{(1\pm i\lambda)/2}(1-z)^{1/2}
     {}_2F_1\!\left(\frac{1\pm i\lambda}{2},
                    \frac{1\pm i\lambda}{2};
                    1\pm i\lambda;z\right).
\]
The ratio, normalized by its leading term at zero, is
\[
 D(z)=\frac{y_+(z)}{y_-(z)}
     =z^{i\lambda}
       \frac{{}_2F_1(a,a;2a;z)}
            {{}_2F_1(\bar a,\bar a;2\bar a;z)}.                \tag{AQ.67.4}
\]
All developing maps with the prescribed Schwarzian differ by a
M\"obius transformation. The ratio in (AQ.67.4) has the same
diagonal monodromy at zero; hence it is a valid developing map
after changing the target half-plane by that transformation.
Its boundary still contains the two fixed points $0,\infty$;
thus it is a straight half-plane bounded by a line through zero,
rather than a disk with a different boundary circle.

We now fix the phase without an unspecified multiple of $\pi$.
On the real interval $0<z<1$, $y_- =\overline{y_+}$.
Their Wronskian, computed at zero, is
$y_+'y_- -y_+y_-'=i\lambda$. In particular neither solution
vanishes there, and the continuously lifted argument $\Theta(z)$
of $D(z)$, normalized by
$\Theta(z)-\lambda\log z\longrightarrow0$ as $z\downarrow0$,
satisfies
\[
 \Theta'(z)=\frac{\lambda}{|y_+(z)|^2}>0.
                                                               \tag{AQ.67.5}
\]
The zero-balanced hypergeometric asymptotic
\[
 {}_2F_1(a,a;2a;z)
   \sim\frac{\Gamma(2a)}{\Gamma(a)^2}[-\log(1-z)]
       \qquad (z\uparrow1)
\]
is DLMF (15.4.21). Thus
\[
 \lim_{z\uparrow1}\Theta(z)=\vartheta(\lambda),\qquad
 \vartheta(\lambda)=2\operatorname{Im}
       \left(\log\Gamma(1+i\lambda)
                    -2\log\Gamma(\tfrac12+i\lambda/2)\right).
                                                               \tag{AQ.67.6}
\]
This is an equality of lifted arguments. For clarity, the
asymptotic is locally uniform in $\lambda\geq0$. The hypergeometric
factor on $0<z<1$ has no zeros when $\lambda>0$ by the Wronskian,
and at $\lambda=0$ it is positive. Its continuous logarithm,
zero at $z=0$, therefore has limiting phase equal to the indicated
continuous gamma phase, rather than that phase plus a nonzero
multiple of $2\pi$. Equivalently that possible integer is locally
constant in $\lambda$ and is zero at $\lambda=0$.

The pair of pants has its reflection fixing the perpendicular
seams. After capping, this antiholomorphic involution fixes
$0,1,\infty$, so it is complex conjugation. The perpendicular
seam from the cusp at $1$ to the geodesic boundary is the terminal
segment of the real interval $(0,1)$. Its developing image lies
on the unit circle, by (AQ.67.4). The cusp endpoint has phase
$\vartheta$. By (AQ.67.5), approaching that cusp from the core
uses phases less than $\vartheta$. Hence the developing half-plane
of the core is the half-plane with argument interval
$(\vartheta-\pi,\vartheta)$. The axis of the boundary monodromy,
whose fixed points are $0,\infty$, is its ray of argument
$\vartheta-\pi/2$. The seam meets that axis perpendicularly.
Starting from the cusp, this is the first intersection of the
seam with the axis; the further formal intersections obtained
after crossing the boundary line of that half-plane are not in
the hyperbolic core. This selects the lifted phase
$\vartheta-\pi/2$, with no integer ambiguity.

Choose $w=1$ at this perpendicular foot. On a collar of the
geodesic, and thence by analytic continuation into its cap,
\[
 w=\exp\left(\frac{\log D-i(\vartheta-\pi/2)}{i\lambda}\right).
\]
Here the logarithm has the lifted phase just specified.
Its boundary modulus is one and its boundary argument has the
normalized arclength, so this is precisely the prescribed cap
coordinate. At zero, (AQ.67.4) gives
\[
 w=\exp\left(-\frac{\vartheta-\pi/2}{\lambda}\right)z+O(z^2),
 \qquad
 r(\ell)=\exp\left(\frac{\vartheta-\pi/2}{\lambda}\right).
                                                               \tag{AQ.67.7}
\]
The duplication formula, DLMF (5.5.5), yields
$\vartheta=2\lambda\log2+2\phi(\lambda/2)$, proving
(AQ.67.1).

It remains to establish the exact range and the nonvanishing
derivative, rather than only an upper bound. The digamma integral
DLMF (5.9.16) gives, for $x\geq0$,
\[
 \phi'(x)=\operatorname{Re}\bigl(\psi(1+ix)
                            -\psi(\tfrac12+ix)\bigr)
           =\int_0^\infty\frac{\cos(xt)}{e^{t/2}+1}\,dt.
                                                               \tag{AQ.67.8}
\]
Let $f(t)=(e^{t/2}+1)^{-1}$. Then
\[
 f''(t)=\frac{e^{t/2}(e^{t/2}-1)}{4(e^{t/2}+1)^3}>0
               \quad(t>0).
\]
The function and its derivatives decay exponentially. Two
integrations by parts, with their endpoint terms included, give
\[
 \phi'(x)=\frac1{x^2}\int_0^\infty f''(t)(1-\cos xt)\,dt>0
                  \quad(x>0).                                \tag{AQ.67.9}
\]
Stirling's expansion for the same analytic logarithms, DLMF
(5.11.1), or the corresponding gamma-ratio expansion (5.11.13),
gives $\phi(x)\longrightarrow\pi/4$ as $x\to\infty$.
Since $\phi(0)=0$, it follows that
$0<\phi(x)<\pi/4$ for $x>0$. Therefore
\[
 \frac{d}{dx}\log r(4\pi x)
       =\frac{x\phi'(x)+\pi/4-\phi(x)}{x^2}>0.                 \tag{AQ.67.10}
\]
At zero, $\phi'(0)=2\log2$, so
\[
 r(\ell)=16\exp\left(-\frac{\pi^2}{\ell}+O(\ell^2)\right)
              \quad(\ell\downarrow0),
 \qquad r(\ell)\longrightarrow4\quad(\ell\to\infty).
\]
Real analyticity and (AQ.67.10) prove that $r$ is the asserted
diffeomorphism.

Finally Fenchel--Nielsen gluing of the two labelled pairs of
pants has exactly one length $\ell>0$ and one twist modulo
$\ell$, equivalently an angle $\theta$ modulo $2\pi$.
The two derivative vectors each have norm $r(\ell)$, and
the twist contributes their tensor's unit phase. This proves
(AQ.67.2). Its inverse is
$\ell=r^{-1}(\sqrt{|q|})$, $\theta=\arg q$ with the latter
viewed as a circle-valued coordinate. In logarithmic polar
coordinates the real Jacobian is diagonal, with entries
$2r'(\ell)/r(\ell)>0$ and $1$. Both the punctured disk and
$\mathbb C^\times$ deformation retract to the same circle
$|q|=1$; the inclusion therefore induces an isomorphism on
$\pi_1$, taking one full twist to the positive generator.
\end{proof}

\subsection{AQ-67-NONHOLO: the natural complex-analytic formulation fails}
Give the labelled geodesic-cut $(0,4)$ space its natural complex
structure, namely the covering
\[
 \mathcal C=\mathcal T(S_{0,4})/\langle D_\alpha\rangle
\]
of the moduli space of four labelled punctures, where $\alpha$ is
the chosen separating curve and $D_\alpha$ is its Dehn twist.
Equivalently, this is the component of the cover that remembers
the isotopy class of the cut. The two sides and all punctures
remain labelled. Fenchel--Nielsen coordinates upstairs are
$(\ell,\tau)\in\mathbb R_{>0}\times\mathbb R$, with
$D_\alpha(\ell,\tau)=(\ell,\tau+\ell)$; downstairs the
normalized angle is $\theta=2\pi\tau/\ell$ modulo $2\pi$.
The canonical capping map of AQ-67-EXACT is not holomorphic
from this complex manifold to $\mathbb C^\times$.
Thus its everywhere invertible real tangent map does not make
it a complex-analytic or $\mathbb C$-algebraic \emph{\'etale}
map. This refutes that stronger interpretation of 67/077--082
for the natural Teichm\"uller complex structure already in
the labelled $(0,4)$ sector.

\begin{proof}
Suppose it were holomorphic. By AQ-67-EXACT it is a bijection
with nonzero real Jacobian onto the punctured disk
$\mathcal D=\{0<|q|<16\}$. The holomorphic inverse function
theorem then makes it a biholomorphism with $\mathcal D$.
Conjugating the rotations $q\mapsto e^{is}q$ by it gives
holomorphic automorphisms of $\mathcal C$. Their lifts to
$\mathcal T(S_{0,4})$, continuously chosen to equal the identity
at $s=0$, are the explicit global diffeomorphisms
\[
 H_s(\ell,\tau)=
       \left(\ell,\tau+\frac{s\ell}{2\pi}\right),
                      \qquad s\in\mathbb R.                  \tag{AQ.67.11}
\]
This follows directly from the tensor's unit twist phase;
any choice of zero-angle section changes $\theta$ by a function
of $\ell$ and leaves (AQ.67.11) unchanged. Since the quotient
projection is a holomorphic covering, every $H_s$ is
holomorphic as well.

Let $\omega_{\mathrm{WP}}$ and $g_{\mathrm{WP}}$ be the
Weil--Petersson K\"ahler form and metric. Wolpert (1983),
Lemma 4.1, p.~229, says that the arclength twist vector field
$t_\alpha=\partial/\partial\tau$ has Hamiltonian potential
$-\ell$, in his normalization:
\[
 \iota_{t_\alpha}\omega_{\mathrm{WP}}=-d\ell.
\]
The vector field of (AQ.67.11) is
$V=(\ell/2\pi)t_\alpha$, so
\[
 \iota_V\omega_{\mathrm{WP}}
   =-\frac{\ell}{2\pi}d\ell
   =-d\!\left(\frac{\ell^2}{4\pi}\right).
\]
Cartan's formula and $d\omega_{\mathrm{WP}}=0$ imply
$\mathcal L_V\omega_{\mathrm{WP}}=0$; hence every global
$H_s$ preserves $\omega_{\mathrm{WP}}$. A different fixed
normalization of the Weil--Petersson form only multiplies
these identities by a nonzero constant and gives the same
conclusion. A holomorphic map preserving the K\"ahler form
preserves the metric, since that metric is recovered from
the form and the complex structure. Thus every $H_s$ would
be a Weil--Petersson isometry of $\mathcal T(S_{0,4})$.

Brock and Margalit, \emph{Weil--Petersson isometries via the
pants complex}, Main Theorem, p.~2 of arXiv:math/0412499v1,
prove that every such isometry is induced by an element of
the extended mapping class group; their theorem explicitly
includes $S_{0,4}$. That group is countable. On the other
hand (AQ.67.11) gives distinct maps for every two distinct
real numbers $s$, because $\ell>0$ and the upstairs twist
coordinate is real, not taken modulo $\ell$. This is an
uncountable family of isometries, a contradiction.
\end{proof}

\subsection{AQ-67-INVOLUTION: the proposition holds; its corollary fails}
The proposition of 67/091 gives an equivalence between
two-regular modules with involution and triples
$(P,Q,m)$, where $P,Q$ are two-regular and
$m\subset P/2P\oplus Q/2Q$ meets each summand trivially.
This proposition is correct, with morphisms of triples
preserving $m$.
\begin{proof}
For an involution on $M$, let $P=M_+$ and $Q=M_-$.
Two-regularity gives $P\cap Q=0$ and
$2M\subset P\oplus Q\subset M$.
Let $m=2M/(2P\oplus2Q)$. If an element of $2M$
has zero $Q$-coordinate, say $2x\in P$, then
$2(\sigma x-x)=0$, so $x\in P$ and $2x\in2P$.
This proves $m\cap(P/2P)=0$, and the other summand
is identical.

Conversely take the inverse image
$L\subset P\oplus Q$ of $m$, with involution
$(p,q)\mapsto(p,-q)$. It is stable because modulo
two this involution is the identity. It is
two-regular, and the intersection conditions give
$L_+=2P\oplus0$, $L_-=0\oplus2Q$.
Multiplication by two canonically identifies these
eigenspaces with $P,Q$ and recovers the original
$m$ from $2L/(4P\oplus4Q)$.
Starting with $M$, the reconstructed $L$ is $2M$,
canonically isomorphic to $M$ by multiplication
by two. Maps restrict to the eigenspaces, and
maps of triples preserve the indicated inverse
images. These constructions are inverse on maps
as well as objects, proving the equivalence.
\end{proof}

The corollary claiming invertible rank-one eigenspaces
for a nontrivial involution on a rank-two projective
module over a commutative two-regular ring is false
even when the ring is an integral domain. Set
\[
 R=\mathbb Z[x,y,z]/(x^2+yz-1),\quad M=R^2,\qquad
 \sigma=\begin{pmatrix}x&y\\z&-x\end{pmatrix}.                   \tag{AQ.9}
\]
\begin{proof}
The polynomial $f=x^2+yz-1$, viewed in
$\mathbb Z[x,y][z]$, is primitive and linear over
the fraction field; hence it is irreducible in
the unique factorization domain $\mathbb Z[x,y,z]$.
Thus $R$ is a domain, in particular two-regular.
Clearly $\sigma^2=1$ and $\sigma\ne\pm1$.

In $S=\mathbb Z[x,y,z]$ set $A=\sigma-I$ and
$B=\sigma+I$. Then $AB=BA=fI$.
If $\bar A\bar v=0$ over $R$, lift $v$ to $S^2$
and write $Av=fw$. Multiplication by $B$ gives
$fv=fBw$; since $f$ is a non-zero-divisor,
$v=Bw$. Therefore $\ker\bar A=\operatorname{im}\bar B$.
Interchanging $A,B$ proves
$\ker\bar B=\operatorname{im}\bar A$, and hence
\[
 P=\ker(\sigma-I)\simeq\operatorname{coker}(\sigma-I).
\]
At the maximal ideal
$\mathfrak m=(2,x-1,y,z)$ every entry of $A$
vanishes, so
$P/\mathfrak mP\simeq\mathbb F_2^2$.
Generically $A$ has rank one, giving $P$ rank one.
An invertible module would have a one-dimensional
fibre at $\mathfrak m$, a contradiction.
The same argument applies to $Q$.
\end{proof}
Connectedness of $\operatorname{Spec}R$ alone cannot
repair the corollary. This algebraic counterexample
does not involve RH, and cannot be excluded by it.

\subsection{AQ-154-COUNT: what can actually be counted over finite fields}
The supplied 154/120 says
\[
 \chi(\mathcal X)=s_0-s_1+s_2=1-g,\qquad
 2s_1=4s_0,\qquad \chi(\mathcal X)=s_2-s_0.
\]
Its $\mathcal X_2$ is a cellular surface of relative
positions of pseudolines. Neither that sheet nor SA11
supplies it with a scheme over a finite field.
The original Euler-characteristic calculation is
therefore not an assertion about
$\#\mathcal X_2(\mathbb F_q)$.

There is a genuine arithmetic consequence for a
different, explicitly specified object.
Let $L_1,\dots,L_n$ be distinct rational projective
lines in $\mathbb P^2_{\mathbb Q}$, let $U$ be their
complement, and let $S$ contain the finitely many
primes where lines coincide or their incidence
lattice changes. If $v$ runs through their
intersection points and $m_v$ lines meet at $v$, set
\[
 b_2=\sum_v(m_v-1)-n+1.
\]
For every $p\notin S$ and $r\geq1$,
\[
 \#U(\mathbb F_{p^r})=p^{2r}-(n-1)p^r+b_2,\qquad
 Z(U_{\mathbb F_p},t)
  =\frac{(1-pt)^{n-1}}{(1-p^2t)(1-t)^{b_2}}.                  \tag{AQ.10}
\]
Consequently its incomplete arithmetic zeta function is
\[
 \zeta^S(U,s)=
 \frac{\zeta^S(s-2)\zeta^S(s)^{b_2}}
      {\zeta^S(s-1)^{n-1}},\qquad
 \zeta^S(s)=\prod_{p\notin S}(1-p^{-s})^{-1}.                 \tag{AQ.11}
\]
\begin{proof}
Every intersection point is rational because the
lines are rational. Over $\mathbb F_{p^r}$ the sum
of the $n$ line counts is $n(p^r+1)$.
It overcounts an intersection $v$ exactly $m_v-1$
times. Subtracting the corrected union count from
$p^{2r}+p^r+1$ proves the first formula.
Insert it in
$Z(U_{\mathbb F_p},t)=
\exp(\sum_{r\geq1}\#U(\mathbb F_{p^r})t^r/r)$
and use $-\log(1-z)=\sum_{r\geq1}z^r/r$ to get
the second. Multiplying its Euler factors gives
(AQ.11), initially for $\operatorname{Re}s>3$
and then by meromorphic continuation.
\end{proof}
Under the adopted RH, the nontrivial-zero contribution
to the divisor of (AQ.11) is on
$\operatorname{Re}s=\tfrac12,\tfrac32,\tfrac52$:
the middle line gives poles when $n>1$, while the
other two give zeros when their exponents are
positive. Trivial zeros, the zeta pole, and the
explicit omitted Euler factors remain separately
specified by (AQ.11). These shifted divisors do not
give a formula for the pseudoline counts or the
topology of $\mathcal X_2$.

\subsection{AQ-SA08-LIFT: a concrete divided-power lifting obstruction}
The supplied SA08 intake concerns the uniqueness problem
of cote 136/117--121 and a proposed lifting input.
The original cote-136 sheets are absent from this packet.
The auxiliary lifting claim can nevertheless be checked
directly in its primary source: Caradot--Lin (2026),
arXiv:2511.22614v4, \S2.1, pp.~5--6, and
Proposition 3.10, pp.~19--20.
There \emph{free complex} means free in each degree,
and surjectivity of the PD ideals is degreewise.
Under these hypotheses the asserted lift need not exist,
over any nonzero commutative ring $k$.
\begin{proof}
Use cochain degrees and the square-zero PD dg algebras
\[
 R'=k\oplus(ku\xrightarrow{\,d\,}kv),\qquad
 |u|=-2,\quad |v|=-1,\quad du=v,
 \qquad R=k\oplus k\epsilon,\quad|\epsilon|=-2,\quad d\epsilon=0.
\]
Their augmentation ideals have zero multiplication,
and $\gamma_n=0$ for $n\geq2$ defines divided powers.
All divided-power identities hold: products of two
ideal elements vanish, and the differential identity
$d\gamma_n(u)=\gamma_{n-1}(u)\,du$ is zero for $n\geq2$,
including $u\,du=uv=0$ when $n=2$.
Define $\pi(u)=\epsilon$, $\pi(v)=0$, and $\pi|_k=1$.
It is a PD dg map surjective on the ideals.
Take the termwise free complex $X=k\epsilon$,
concentrated in degree $-2$, and the identity chain
map $X\to k\epsilon\subset R$.
A PD dg algebra lift from the free PD dg algebra
on $X$ would restrict on its generator to a chain
lift $\epsilon\mapsto cu$.
Compatibility with $\pi$ forces $c=1$, whereas
compatibility with the differential forces $cv=0$,
hence $c=0$ in $k$. This contradicts $k\ne0$.
\end{proof}
The notation $X$ here denotes just $k$ placed in
degree $-2$; no resolution or projectivity in the
category of complexes is being assumed.
The failed implication is precisely from a graded
lift to a chain lift. Theorem 3.15 of the same
version, pp.~23--24, assumes a quasi-isomorphism.
Our $\pi$ is not one:
$H^*(R')=k$ but $H^*(R)=k\oplus k\epsilon$.
Consequently this counterexample does not refute
that later theorem.

The unresolved SA08 statements are the actual
integral common-source assertion A-CS, terminality
of the all-arrow localization A-TL, and application
of the SA09 comparison structures A-SA09.
The intake states them but does not supply their
complete definitions and 15,000-line proof context.
Neither a weight grading on a divided-power algebra
nor the adopted analytic zero theorem supplies a
chain lift in the displayed counterexample.
That is a specific false lifting route; it is not
a disproof of the remaining existence or
localization conjectures.


